%% LyX 2.4.4 created this file.  For more info, see https://www.lyx.org/.
%% Do not edit unless you really know what you are doing.
\documentclass[oneside,english]{book}
\usepackage[T1]{fontenc}
\usepackage[latin9]{inputenc}
\setcounter{secnumdepth}{3}
\setcounter{tocdepth}{3}
\usepackage{units}
\usepackage{amsmath}
\usepackage{amsthm}
\usepackage{amssymb}
\usepackage{cancel}
\usepackage{geometry}
\geometry{verbose,tmargin=1in,bmargin=1in,lmargin=1.5in,rmargin=1.5in}
\usepackage{wasysym}
\PassOptionsToPackage{normalem}{ulem}
\usepackage{ulem}

\makeatletter

%%%%%%%%%%%%%%%%%%%%%%%%%%%%%% LyX specific LaTeX commands.
%% Because html converters don't know tabularnewline
\providecommand{\tabularnewline}{\\}

%%%%%%%%%%%%%%%%%%%%%%%%%%%%%% Textclass specific LaTeX commands.
\theoremstyle{definition}
\ifx\thechapter\undefined
  \newtheorem{example}{\protect\examplename}
\else
  \newtheorem{example}{\protect\examplename}[chapter]
\fi
\theoremstyle{definition}
\ifx\thechapter\undefined
  \newtheorem{defn}{\protect\definitionname}
\else
  \newtheorem{defn}{\protect\definitionname}[chapter]
\fi
\theoremstyle{definition}
\ifx\thechapter\undefined
  \newtheorem{xca}{\protect\exercisename}
\else
  \newtheorem{xca}{\protect\exercisename}[chapter]
\fi
\theoremstyle{remark}
\ifx\thechapter\undefined
  \newtheorem{rem}{\protect\remarkname}
\else
  \newtheorem{rem}{\protect\remarkname}[chapter]
\fi
\theoremstyle{plain}
\ifx\thechapter\undefined
  \newtheorem{prop}{\protect\propositionname}
\else
  \newtheorem{prop}{\protect\propositionname}[chapter]
\fi
\ifx\proof\undefined\
  \newenvironment{proof}[1][\proofname]{\par
    \normalfont\topsep6\p@\@plus6\p@\relax
    \trivlist
    \itemindent\parindent
    \item[\hskip\labelsep
          \scshape
      #1]\ignorespaces
  }{%
    \endtrivlist\@endpefalse
  }
  \providecommand{\proofname}{Proof}
\fi
\theoremstyle{plain}
\ifx\thechapter\undefined
  \newtheorem{lem}{\protect\lemmaname}
\else
  \newtheorem{lem}{\protect\lemmaname}[chapter]
\fi
\theoremstyle{plain}
\ifx\thechapter\undefined
  \newtheorem{thm}{\protect\theoremname}
\else
  \newtheorem{thm}{\protect\theoremname}[chapter]
\fi
\theoremstyle{plain}
\ifx\thechapter\undefined
  \newtheorem{cor}{\protect\corollaryname}
\else
  \newtheorem{cor}{\protect\corollaryname}[chapter]
\fi

%%%%%%%%%%%%%%%%%%%%%%%%%%%%%% User specified LaTeX commands.
\usepackage{hyperref}
\usepackage[usenames,dvipsnames]{xcolor} %adds additional color options
%\hypersetup{colorlinks=true, urlcolor=Blue}%Blue
\hypersetup{colorlinks=true, urlcolor=Blue,linkcolor=[rgb]{0,0,.5}} %blue

\usepackage{multicol}

\usepackage{tikz}
\usepackage{tikz-3dplot}

\usetikzlibrary{calc,arrows,shapes,backgrounds,matrix,shapes.geometric,fit,trees,positioning}

\usetikzlibrary{patterns}
\usetikzlibrary{plotmarks}
\usepackage{pgfplots}
\pgfplotsset{compat=newest}

\usepackage{calc}
\usepackage[nomessages]{fp}

\usepackage{datetime}
%\usepackage{subcaption}

%\usepackage{amsthm}

%%%%%commands
\newcommand{\mb}[1]{$\mathbb{#1}$}
\newcommand{\funct}[2]{$f\colon#1\rightarrow#2$}
% Added by lyx2lyx
\setlength{\parskip}{\smallskipamount}
\setlength{\parindent}{0pt}

\makeatother

\usepackage{babel}
\providecommand{\corollaryname}{Corollary}
\providecommand{\definitionname}{Definition}
\providecommand{\examplename}{Example}
\providecommand{\exercisename}{Exercise}
\providecommand{\lemmaname}{Lemma}
\providecommand{\propositionname}{Proposition}
\providecommand{\remarkname}{Remark}
\providecommand{\theoremname}{Theorem}

\begin{document}
\title{MTH 335: Introduction to Abstract Algebra Notes}

\maketitle
\tableofcontents{}

%% LyX 2.4.4 created this file.  For more info, see https://www.lyx.org/.
%% Do not edit unless you really know what you are doing.
\documentclass[oneside,english]{book}
\usepackage[T1]{fontenc}
\usepackage[latin9]{inputenc}
\setcounter{secnumdepth}{3}
\setcounter{tocdepth}{3}
\usepackage{units}
\usepackage{amsthm}
\usepackage{amssymb}
\usepackage{geometry}
\geometry{verbose,tmargin=1cm,bmargin=1cm,lmargin=1.5cm,rmargin=1.5cm}
\PassOptionsToPackage{normalem}{ulem}
\usepackage{ulem}

\makeatletter
%%%%%%%%%%%%%%%%%%%%%%%%%%%%%% Textclass specific LaTeX commands.
\theoremstyle{definition}
\newtheorem{example}{\protect\examplename}
\theoremstyle{definition}
\newtheorem{defn}{\protect\definitionname}
\theoremstyle{definition}
\newtheorem{xca}{\protect\exercisename}
\theoremstyle{remark}
\newtheorem{rem}{\protect\remarkname}
\theoremstyle{plain}
\newtheorem{prop}{\protect\propositionname}
\ifx\proof\undefined\
  \newenvironment{proof}[1][\proofname]{\par
    \normalfont\topsep6\p@\@plus6\p@\relax
    \trivlist
    \itemindent\parindent
    \item[\hskip\labelsep
          \scshape
      #1]\ignorespaces
  }{%
    \endtrivlist\@endpefalse
  }
  \providecommand{\proofname}{Proof}
\fi
\theoremstyle{plain}
\newtheorem{lem}{\protect\lemmaname}
\theoremstyle{plain}
\newtheorem{thm}{\protect\theoremname}

%%%%%%%%%%%%%%%%%%%%%%%%%%%%%% User specified LaTeX commands.
\usepackage{hyperref}
\usepackage[usenames,dvipsnames]{xcolor} %adds additional color options
\hypersetup{colorlinks=true, urlcolor=Blue}%Blue

\usepackage{multicol}

\usepackage{tikz}
\usepackage{tikz-3dplot}
\usetikzlibrary{calc,arrows,shapes,backgrounds,matrix,shapes.geometric,fit,trees,positioning}

\usetikzlibrary{patterns}
\usetikzlibrary{plotmarks}
\usepackage{pgfplots}
\pgfplotsset{compat=newest}

\usepackage{calc}
\usepackage[nomessages]{fp}

\usepackage{datetime}

%%%%%commands
\newcommand{\mb}[1]{$\mathbb{#1}$}
\newcommand{\funct}[2]{$f\colon#1\rightarrow#2$}
% Added by lyx2lyx
\setlength{\parskip}{\smallskipamount}
\setlength{\parindent}{0pt}

\makeatother

\usepackage{babel}
\providecommand{\definitionname}{Definition}
\providecommand{\examplename}{Example}
\providecommand{\exercisename}{Exercise}
\providecommand{\lemmaname}{Lemma}
\providecommand{\propositionname}{Proposition}
\providecommand{\remarkname}{Remark}
\providecommand{\theoremname}{Theorem}

\begin{document}

\chapter{Introduction to Sets}

Set theory is a foundation on which (almost) all of mathematics can
be built. 
\begin{itemize}
\item A \emph{set }is a collection of distinct objects called \emph{elements}.
\end{itemize}
This is not really a definition (for a rigorous definition see this
\href{http://en.wikipedia.org/wiki/Set_theory}{link}); to avoid some
messiness we will use this ``intuitive'' (sometimes called ``naive'')
understanding of set and element. Which has good enough for Cantor,
so it should be good enough for us.\footnote{George Cantor founded set theory using the ``naive'' defintions
in 1874 with a single paper. However by 1900 a number of contradictions
were found. Most famously Russell's paradox: ``the set of all sets
that are not members of themselves.'' This set must be a member of
itself, and yet is not a member of itself!} 

\section{Describing Sets}

Lots of new notation will be introduced as we move through the material.
The trade-off for learning all the new notation is certainly worth
it. Think of trying to solve even a simple quadratic, $x^{2}+2x-10=0$,
without the advantage of using $x$ as the unknown and Arabic numbers.
\begin{itemize}
\item Sets are usually denoted with capital letters, e.g., $A,B,C,\dots$
etc; and elements in sets are usually denoted with lowercase letters
$a,b,c,\dots$.
\item If $a$ is an element in the set $A$, we write $a\in A$. If $a$
is not in $A$ we write $a\notin A$ (slashes usually mean negation). 
\end{itemize}
Sets are described in a number of different ways:

Described: 
\begin{example}
$a,b,c$ are elements of the set $A$; also written $a,b,c\in A$.
\end{example}

\begin{example}
$\mathbb{R}$ is the set of all real numbers.
\end{example}

\emph{Tabular }form: 
\begin{example}
$A=\left\{ a,b,c\right\} $, means $a,b,c\in A$.

\begin{itemize}
\item []Dots are used when unlisted elements can be implied:
\end{itemize}
\end{example}

\begin{example}
$A=\left\{ a,b,c,\dots z\right\} $, it is implied that all the letters
in alphabet are in $A$.
\end{example}

\begin{example}
$\mathbb{Z}=\left\{ \dots,-1,0,1,2,\dots\right\} $
\end{example}
\emph{Set builder }form: 
\begin{example}
$\mathbb{Z}^{+}=\left\{ n\in\mathbb{Z}\colon\ensuremath{n>0}\right\} $,
read ``$n$ is an integer \emph{such that }$n$ is greater than $0$.''
\end{example}

\begin{example}
$2\mathbb{Z}=\left\{ n\in\mathbb{Z}\colon n=2m\mbox{ for some }m\in\mathbb{Z}\right\} $.
This set is all the even numbers. So $4\in2\mathbb{Z}$ and $11\notin2\mathbb{Z}$.
\end{example}

\section{Some Basic Definitions}

\emph{A farmer had a 100 feet of fence and wanted to enclose as large
an area as possible. After a little thought, an engineer made a square
resulting in an area of $25^{2}=625\mbox{ ft}^{2}$. A physicist (knowing
he could do better) made a circle resulting in $\pi\left(\nicefrac{100}{2\pi}\right)^{2}\approx795.7\,\mbox{ft}^{2}$.
Finally a mathematician steps forward, breaks off a small piece of
fence, and says ``I define myself to be on the outside!'' (resulting
in an area approximately equal to the area of the Earth). }

How we define things is very important. When trying to prove something,
read the definitions very carefully. For the following definitions,
assume that $A$ and $B$ are both sets. 
\begin{defn}
$A$ is a \emph{subset of $B$ }if every element of $A$ is also an
element of $B$, written $A\subseteq B$. 
\end{defn}
\begin{example}
If $A=\{1,2,3\}$ and $B=\{0,1,2,3\}$then $A\subseteq B$ and $B\not\subseteq A$.
Also note that $A\subseteq A$. 
\end{example}
\begin{itemize}
\item $A\subseteq A$ for any set $A$. Every set is subset of itself. Since
every element of a set $A$ will also be an element of $A$.
\end{itemize}
\begin{example}
$\mathbb{N}\subseteq\mathbb{Z}\subseteq\mathbb{Q}\mathbb{\subseteq R\subseteq C}$:
The natural, integers, rational, real, and complex numbers respectively.
\end{example}
\begin{defn}
The set with no elements is called the \emph{empty set}, denoted $\emptyset$.
\end{defn}
\begin{itemize}
\item $\emptyset=\{x\in A\colon x\notin A\mbox{ for any set }A\}$.
\item Let $B$ be set of numbers that are both irrational and rational real
numbers. $B=\emptyset$.
\item $\emptyset\subseteq A$ for any set $A$, i.e., \emph{the empty set
is a subset of every set}. 
\end{itemize}
\begin{defn}
$A=B$ if and only if $A$ and $B$ have the same elements.
\end{defn}
\begin{itemize}
\item It is often helpful to use the equivalent definition: $A=B$ is $A\subseteq B$
\uline{and} $B\subseteq A$.
\end{itemize}
\begin{example}
Let $A=\{a\colon a=2^{n}\mbox{ for any integer }n\}$ and $B=\{b\colon b=2^{n-2}$
for any integer $n$\}. Show that $A=B$. (Show that $A\subseteq B$)
Let $a\in A$. If $n$ is an integer then there is another integer
$m$ such that $m=n-2$. So then $a=2^{n}=2^{m-2}\in B$. \\
(Show that $B\subseteq A$) Let $b\in B$. If $n$ is an integer then
there is another integer $m$ such that $n-2=m$. So then $b=2^{n-2}=2^{m}\in A$.

\begin{itemize}
\item Note that proving $B\subseteq A$ was nearly identical to proving
$A\subseteq B$ (I literally copied,pasted, and switch a few letters).
When one part of prove is so clearly similar to an already completed
previous part, we will often just write ``similarly it can be shown
that $B\subseteq A$''.
\end{itemize}
\end{example}
\begin{defn}
$A$ is a \emph{proper subset of $B$ }if every element of $A$ is
also an element of $B$ \uline{and} $A\neq B$, written $A\subseteq B$.
\end{defn}
\begin{example}
Let $A=\{a\in\mathbb{Z}\colon a=2^{n}\mbox{ for a positive integer \ensuremath{n}}\}$
and $B=\{a\in\mathbb{Z}\colon a=2n\}$. Show that $A$ is a proper
subset of $B$. 

\begin{itemize}
\item []We must show that $A\subseteq B$ but $A\neq B$. Let $a\in A$.
by definition $a=2^{n}=2\cdot2^{n-1}$ for $n\geq0$. Since $2^{n-1}$
is an integer, $a\in A$. (To show that $A\neq B$ we need only find
1 element in $B$ that is not in $A$. There are many.) $6=2\cdot3\in B$
but $6\notin A$ since $6\neq2^{n}$ for $n$ a positive integer. 
\end{itemize}
\end{example}
\begin{itemize}
\item Our use of$\subseteq$ and $\subset$ is analogous to $\leq$ and
$<$.\footnote{Note that the meaning of set notation varies among authors. $\varsubsetneq$
is often used to denote proper subset, and others use$\subset$ for
subset. }
\end{itemize}
\begin{figure}
\hfill%
\begin{minipage}[t]{0.33\textwidth}%
\begin{center}

\begin{tikzpicture}     	\begin{scope}[scale=.5,shift={(3cm,-5cm)}, fill opacity=0.5,mytext/.style={text opacity=1,font=\large\bfseries}]         
		\draw[fill=red, draw = black] (0,0) circle (2);         
		\draw[fill=green, draw = black] (-0.75,0) circle (1);     
		\node[mytext] at (0,1.5) (B) {\large\textbf{B}};     		
		\node[mytext] at (-1,0) (A) {\large\textbf{A}};     	
	\end{scope}
\end{tikzpicture}
\caption{}{$A\subseteq B$ and $A\subset B$}
\par\end{center}%
\end{minipage}\hfill%
\begin{minipage}[t]{0.33\textwidth}%
\begin{center}

\begin{tikzpicture}     	\begin{scope}[scale=.5,shift={(3cm,-5cm)}, fill opacity=0.5,mytext/.style={text opacity=1,font=\large\bfseries}]         
		\draw[fill=red, draw = black] (0,0) circle (2);         
		\draw[fill=green, draw = black] (0.0,0) circle (2);     
		\node[mytext] at (0,1.5) (B) {\large\textbf{B}};     		
		\node[mytext] at (-1,0) (A) {\large\textbf{A}};     	
	\end{scope}
\end{tikzpicture}
\caption{}{$A\subseteq B$ and $A\not\subset B$}
\par\end{center}%
\end{minipage}\hfill%
\begin{minipage}[t]{0.33\textwidth}%
\begin{center}

\begin{tikzpicture}     	\begin{scope}[scale=.5,shift={(3cm,-5cm)}, fill opacity=0.5,mytext/.style={text opacity=1,font=\large\bfseries}]         
		\draw[fill=red, draw = black] (0,0) circle (2);         
		\draw[fill=green, draw = black] (-2.0,0) circle (1);     
		\node[mytext] at (0,1.5) (A) {\large\textbf{B}};     		
		\node[mytext] at (-2.5,0) (B) {\large\textbf{A}};     	
	\end{scope}
\end{tikzpicture}
\caption{}{$A\not\subset B$}
\par\end{center}%
\end{minipage}\hfill

\caption{Venn diagrams}
\end{figure}

\begin{defn}
Let $A$ be any set. The set of all subsets of $A$ is called the
\emph{power set }of $A$, denoted $\mathcal{P}(A)$. 
\end{defn}
\begin{xca}
Let $A=\{0,1,2,3\}$. Find $\mathcal{P}(A)$ (don't forget $A$ and
$\emptyset$).
\end{xca}
\begin{itemize}
\item Note that: $A\in\mathcal{P}(A)$, $\emptyset\in\mathcal{P}(A)$, and
$\mathcal{P}(\emptyset)=\emptyset$.
\item \label{thm:order-of-power-set}Theorem \ref{thm:-power set order}:
If $A$ be any set with $n$ elements then $\mathcal{P}(A)$ has $2^{n}$
elements. 
\end{itemize}
\begin{rem}
Some summary.

\begin{enumerate}
\item $A=B\iff A\subseteq B$ and $B\subseteq A$.\footnote{``$\iff$'' is equivalent to writing ``if and only if''. Often
this is notated ``iff.''}
\item $A\subseteq A$
\item $\emptyset\subseteq A$
\item $A\in\mathcal{P}(A)$ and $\emptyset\in\mathcal{P}(A)$
\end{enumerate}
\end{rem}

\section{Set Operations}
\begin{defn}
The \emph{union }of sets $A$ and $B$ is defined by $A\cup B=\{x\colon x\in A\mbox{ or }x\in B\}$
\end{defn}
\begin{example}
Let $A=\{1,2,3\}$ and $B=\{2,3,4\}$. $A\cup B=\{1,2,3,4\}$. 
\end{example}
\begin{xca}
$A\cup B=B$ if and only if $A\subseteq B$.
\end{xca}

\begin{example}
$\mathbb{R\cup\mathbb{Z}=\mathbb{R}}$. $\mathbb{Z}\subseteq\mathbb{R}$
so this is true by the above result.
\end{example}

\begin{xca}
$A\cup B=B\cup A$.
\end{xca}
\begin{defn}
The intersection of sets $A$ and $B$ is defined by $A\cap B=\{x\colon x\in A\mbox{ and }x\in B\}$.
\end{defn}
\begin{example}
Let $A=\{1,2,3\}$ and $B=\{2,3,4\}$. $A\cap B=\{2,3\}$.
\end{example}

\begin{example}
Let $A=\{1,2,3\}$ and $B=\{4,5,6\}$. $A\cap B=\emptyset$.
\end{example}

\begin{xca}
$A\cap B=B$ if and only if $B\subseteq A$.
\end{xca}

\begin{xca}
$A\cap B=B\cap A$.
\end{xca}
\begin{figure}[h]
\begin{centering}
\def\firstcircle{(90:1.75cm) circle (2.5cm)} 
\def\secondcircle{(210:1.75cm) circle (2.5cm)} 
\def\thirdcircle{(330:1.75cm) circle (2.5cm)}

\begin{tikzpicture}     
\begin{scope}     
	\fill[red]\firstcircle;       
	\fill[green] \secondcircle;       
	\fill[blue] \thirdcircle;     
\end{scope}     
\begin{scope}
	\clip 
	\secondcircle;       
	\fill[yellow] 
	\firstcircle;     
\end{scope}     
\begin{scope}       
	\clip 
	\secondcircle;       
	\fill[cyan] 
	\thirdcircle;     
\end{scope}     
\begin{scope}       
	\clip 
	\firstcircle;       
	\fill[magenta] 
	\thirdcircle;     
\end{scope}     
\begin{scope}       
	\clip \firstcircle;       
	\clip \secondcircle;       
	\fill[white] \thirdcircle;     
\end{scope}     
\draw 
	\firstcircle node[text=white,above] {$A$};     
	\draw \secondcircle node [text=white,below left] {$B$};     
	\draw \thirdcircle node [text=white,below right] {$C$};     
\begin{scope}[font=\small]       
	\node(0,0) {$A\cap B\cap C$};       
	\draw(30:1.5cm) node {$A\cap C$};       
	\draw(150:1.5cm) node {$A\cap B$};       
\draw(270:1.5cm) node {$B\cap C$};     
\end{scope}   
\end{tikzpicture}\caption{A Venn diagram showing the intersections of three sets, $A$, $B$,
and$C$. }
\par\end{centering}
\end{figure}

\begin{defn}
The \emph{difference of sets }$B$ and $A$ is the set of elements
that are in $B$ and not in $A$; written $B-A$.\footnote{Also called the \emph{relative complement of }$A$ in $B$ usually

denoted $B\setminus A$.}
\end{defn}
\begin{itemize}
\item $B-A=\{x\colon x\in B\mbox{ but }x\notin A\}$ (set builder form).
\end{itemize}
\begin{example}
$\mathbb{Z}-2\mathbb{Z}$ is the set of odd integers.
\end{example}
\begin{defn}
The compliment of a set $A$ is the set of elements (under consideration)
not in $A$; denoted $A^{c}$.\footnote{We are being a little sloppy here by ignoring some complications of

the universal set in trade for an easier introduction to sets. }
\end{defn}
\begin{itemize}
\item $A^{c}=\{x\colon x\notin A\}$.
\end{itemize}
\begin{example}
When considering real numbers, $\mathbb{Z}^{c}=\{x\colon x\notin\mathbb{Z}\}$
is the set of all numbers that are not integers. 
\end{example}
\begin{prop}
$A-B=A\cap B^{c}$.

\begin{proof}
Let $x\in A-B$ so then $x\in A$ and $x\notin B$. By definition
of compliments $x\in B^{c}$. So then $x\in A-B$ and $A-B\subseteq A\cap B^{c}$.

Let $x\in A\cap B^{c}$ so then $x\in A$ and $x\notin B$. By definition
of difference of sets, $x\in A-B$ thus $A\cap B^{c}\subseteq A-B$.

$A-B\subseteq A\cap B^{c}$ and $A\cap B^{c}\subseteq A-B$ imply
$A\cap B^{c}=A-B$.
\end{proof}
\end{prop}
\begin{defn}
The number of elements in a set $A$ is called the \emph{order of
$A$}; denoted $|A|$.\footnote{Also called the \emph{cardinality of $A$.}}
\end{defn}
\begin{itemize}
\item When we write $|A|=n$ it is assumed that $n$ is a non-negative integer.
\end{itemize}
\begin{example}
Let $A=\{a,b,c\}$. $|A|=3$.
\end{example}

\begin{example}
Let $A=\{a,b,c\}$ then $|\mathcal{P}(A)|=2^{3}=8$ by Theorem \ref{thm:order-of-power-set}.
\end{example}

\begin{example}
$|\emptyset|=0$.
\end{example}
\begin{lem}
\label{lem:lemma-for-power-setr-num}If $|A|=n$ and $|B|=n+1$ for
sets $A$ and $B$ then $|\mathcal{P}(B)|=2|\mathcal{P}(A)|$.

\begin{proof}
We can without loss of generality\footnote{Often abbreviated WLOG, ``without loss of generality'' is used before
an assumption in a proof to narrow the premise to some special case
when the proof using the special case can obviously be adapted for
all other cases.} assume that $B=A\cup\{b\}$. So that $A=B-\{a\}$. $B$ has two kinds
of subsets; those containing $b$ and those that do not. Now let $S\in\mathcal{P}(B)$
such that $b\notin S$. Then $S\in\mathcal{P}(A)$, and there are
$n$ such sets. Now let $S'\in\mathcal{P}(B)$ such that $b\in S'$.
Then $S'-\{b\}\subseteq A$, and there are $n$ such sets. As a subsets
of $B$ either contains $b$ or it does not, $|\mathcal{P}(B)|=|\mathcal{P}(A)|+|\mathcal{P}(A)|=2|\mathcal{P}(A)|$.
\end{proof}
\end{lem}
\begin{thm}
\textup{\label{thm:-power set order}$|\mathcal{P}(A)|=2^{n}$ for
any set $A$.}
\end{thm}
\begin{example}
Prove theorem \ref{thm:order-of-power-set}. 
\end{example}
\begin{itemize}
\item []Hint: Use lemma \ref{lem:lemma-for-power-setr-num} and induction.
\end{itemize}

\section{Exercises}
\begin{enumerate}
\item Write in tabular form: 

\begin{enumerate}
\item $A=\{x\colon x\mbox{ is a letter in "do not repeat"}\}$

\begin{enumerate}
\item $3\mathbb{Z}=\left\{ x\in\mathbb{Z}\colon x=3n\mbox{ for some \ensuremath{n\in\mathbb{Z}}}\right\} $
\item $N=\left\{ x\colon x\mbox{ is odd and \ensuremath{x}is even}\right\} $
(recall that $0$ is even)
\end{enumerate}
\item Write these sets in set-builder form:

\begin{enumerate}
\item Let $A=\{3,4,5,6,7,\dots\}$
\item Let $B=\{3,5,7,9,\dots\}$
\item Let $\mathbb{Q}$ be the set of rational numbers. 
\end{enumerate}
\end{enumerate}
\item Find all sets of $A=\{1,2,3\}$.
\item Prove that every nonempty subset has a proper subset. 
\item Prove that if $A$ is a subset of $\emptyset$, then $A=\emptyset$.
\item Prove that $A=\{2,3,4,5\}$ is not a subset of $B=\{x\mid x\mbox{ is even}\}$.
\item Prove that $A$ and $B$ are subsets of $A\cup B$.
\item Prove that $A=A\cup A$ and $A=A\cap A$.
\item Prove that $A\cap B$ is a subset of $A$ and of $B$.
\item Prove that $A\cup B=\emptyset$ implies that $A=\emptyset$ and $B=\emptyset$.
\item Prove that $A\cup B=B$ if and only if $A\subseteq B$.
\item Prove that $A\cap B=B$ if and only if $B\subseteq A$.
\item Prove that $A\cup\left(A\cap B\right)=A$.
\item Prove that $A\cap\left(A\cup B\right)=A$.
\item Prove that $\mathcal{P}(A\cap B)=\mathcal{P}(A)\cap\mathcal{P}(B)$.
\item Prove that $A-B$ is a subset of $A\cup B$.
\item Prove that if $A$ a subset of $B$ then $A\cap B=A$.
\item Let $A\cap B=\emptyset$. Prove that $B\cap A'=B$.
\item Prove that if $A$ is a subset of $B$ then $A\cup B=B$.
\item Prove that $A'-B'=B-A$.
\item Prove that if $A$ is a subset of $B$ then $B'$ is a subset of $A'$.
\item Prove that if $A\cap B=\emptyset$ then $A\cup B'=B'$.
\item Prove that $\left(A\cap B\right)'=A'\cup B'$.
\item Prove that $A\subset B$ implies that $A\cup\left(B-A\right)=B$.
\item $\bigstar$ The \emph{product set }of sets $A$ and $B$ is the set
of all ordered pairs $(a,b)$ where $a\in A$ and $b\in B$; it is
denoted by $A\times B$. Prove that $A\times(B\cap C)=(A\times B)\cap(A\times C)$. 
\end{enumerate}

\end{document}


%% LyX 2.4.4 created this file.  For more info, see https://www.lyx.org/.
%% Do not edit unless you really know what you are doing.
\documentclass[oneside,english]{book}
\usepackage[T1]{fontenc}
\usepackage[latin9]{inputenc}
\setcounter{secnumdepth}{3}
\setcounter{tocdepth}{3}
\usepackage{units}
\usepackage{amstext}
\usepackage{amsthm}
\usepackage{amssymb}
\usepackage{cancel}
\usepackage{geometry}
\geometry{verbose,tmargin=1cm,bmargin=1cm,lmargin=1.5cm,rmargin=1.5cm}
\usepackage{wasysym}
\PassOptionsToPackage{normalem}{ulem}
\usepackage{ulem}

\makeatletter

%%%%%%%%%%%%%%%%%%%%%%%%%%%%%% LyX specific LaTeX commands.
%% Because html converters don't know tabularnewline
\providecommand{\tabularnewline}{\\}

%%%%%%%%%%%%%%%%%%%%%%%%%%%%%% Textclass specific LaTeX commands.
\theoremstyle{definition}
\newtheorem{example}{\protect\examplename}
\theoremstyle{definition}
\newtheorem{defn}{\protect\definitionname}
\ifx\proof\undefined\
  \newenvironment{proof}[1][\proofname]{\par
    \normalfont\topsep6\p@\@plus6\p@\relax
    \trivlist
    \itemindent\parindent
    \item[\hskip\labelsep
          \scshape
      #1]\ignorespaces
  }{%
    \endtrivlist\@endpefalse
  }
  \providecommand{\proofname}{Proof}
\fi
\theoremstyle{definition}
\newtheorem{xca}{\protect\exercisename}
\theoremstyle{plain}
\newtheorem{prop}{\protect\propositionname}
\theoremstyle{plain}
\newtheorem{thm}{\protect\theoremname}

%%%%%%%%%%%%%%%%%%%%%%%%%%%%%% User specified LaTeX commands.
\usepackage{hyperref}
\usepackage[usenames,dvipsnames]{xcolor} %adds additional color options
\hypersetup{colorlinks=true, urlcolor=Blue}%Blue

\usepackage{multicol}

\usepackage{tikz}
\usepackage{tikz-3dplot}
\usetikzlibrary{calc,arrows,shapes,backgrounds,matrix,shapes.geometric,fit,trees,positioning}

\usetikzlibrary{patterns}
\usetikzlibrary{plotmarks}
\usepackage{pgfplots}
\pgfplotsset{compat=newest}

\usepackage{calc}
\usepackage[nomessages]{fp}

\usepackage{datetime}

%%%%%commands
\newcommand{\mb}[1]{$\mathbb{#1}$}
\newcommand{\funct}[2]{$f\colon#1\rightarrow#2$}
% Added by lyx2lyx
\setlength{\parskip}{\smallskipamount}
\setlength{\parindent}{0pt}

\makeatother

\usepackage{babel}
\providecommand{\definitionname}{Definition}
\providecommand{\examplename}{Example}
\providecommand{\exercisename}{Exercise}
\providecommand{\propositionname}{Proposition}
\providecommand{\theoremname}{Theorem}

\begin{document}

\chapter{Functions and Relations}

\section{Functions}

The previous chapter covered sets and their basic operations. Now
we introduce the notion of a function between sets. Just as in calculus,
a function will assign a unique output for every input. Familiar notations
such as, $f$ and $f(x)$ will be used. The same concepts will be
used; only the inputs and outputs will include things other than numbers.
A \emph{function }(or \emph{mapping) }from sets $A$ to $B$\emph{,
$f\colon A\rightarrow B$,} is a rule that assigns each element in
$A$ to exactly one element in $B$.
\begin{itemize}
\item $A$ is called the \emph{domain }of $f$. 

\begin{itemize}
\item $B$ is called the \emph{range }of $f$.\footnote{Also called the \emph{co-domain.}}
\item If $f(a)=b$ for some $a\in A$, then $b$ is called the\emph{ image
of $a$ under $f$.}
\item $\{b\colon b\in B\mbox{ and }f(a)=b\mbox{ for some }a\in A\}$ is
called the \emph{image of} $A$ \emph{under $f$}.
\end{itemize}
\item Functions are well defined. Meaning that an input is assigned to one
and only output. 
\end{itemize}
\begin{example}
Let $f\colon\mathbb{R}\rightarrow\mathbb{R}$. Then $f(x)=x^{2}+1$
is a function. it is clear from the equation that $f$ is well defined.
\end{example}

\begin{example}
Let $f\colon\mathbb{R}\rightarrow\mathbb{R}$. Then $f(x)=\pm\sqrt{x}$
is \uline{not} a function. Since $f(1)=+1$ and $-1$.
\end{example}

\begin{example}
\label{ex-floor-funct}Consider $f\colon\mathbb{R}\rightarrow\mathbb{Z}$
given by the \emph{floor function $f(x)=\left\lfloor x\right\rfloor $
}where $\left\lfloor x\right\rfloor $ is the largest integer less
than or equal to $x$, e.g., $\left\lfloor \pi\right\rfloor =3$.
Note that the floor function is also a function from $\mathbb{R}$
to $\mathbb{R}$, $f\colon\mathbb{R}\rightarrow\mathbb{Z}$.
\end{example}

\begin{example}
Let $f\colon\mathbb{Z}\times\mathbb{Z}\rightarrow\mathbb{Z}$ given
by $f(x,y)\rightarrow xy$. Then $f$ is a function. 
\end{example}

\begin{example}
Let $A=\{a,b,c,\dots,z\}$ and $f\colon A\rightarrow\mathbb{Z}$ be
function defined by assigning the $n^{\mbox{th}}$letter to $n$,
e.g., $f(3)=c$.
\end{example}
\begin{defn}
Let $f\colon B\rightarrow C$ and $g\colon A\rightarrow B$ be functions,
then the \emph{composition }of $f$ and $g$ is $f\circ g\colon A\rightarrow C$
defined as $f(g(a))=c$ where $g(a)=b$ and $f(b)=c$ for elements
$a\in A$, $b\in B$, and $c\in C$.
\end{defn}
\begin{figure}
\begin{centering}
\begin{tikzpicture}[description/.style={fill=white,inner sep=2pt}]   
	\matrix (m)[matrix of math nodes,     
	row sep=3em,column sep=2.5em,     
	text height=1.5ex, text depth=0.25ex,     	nodes={ellipse,draw,minimum width=2cm},   	]
	{
		A && B \\     
		& C & 
		\\   };   
	\path[->,font=\scriptsize]   (m-1-1) edge node[auto] {$ f $} (m-1-3)   edge node[description] {$ g \circ f $} (m-2-2)   (m-1-3) edge node[auto] {$ g $} (m-2-2); \end{tikzpicture}
\par\end{centering}
\caption{Composition of functions $f$ and $g$.}
\end{figure}

\begin{itemize}
\item In general the composition of function is not commutative, i.e., $f\circ g\neq g\circ f$.
\end{itemize}
\begin{example}
Let $f,g:\mathbb{R}\rightarrow\mathbb{R}$ and $f(x)=2x+1$ and $g(x)=x^{2}$.
Then $f\circ g(x)=f(g(x))=2(x^{2})+1=2x^{2}+1$.
\end{example}

\begin{defn}
A function, \funct{A}{B}, is called \emph{one-to-one}\footnote{Also called an \emph{injective function}}
if $f(a_{1})=f(a_{2})$ implies that $a_{1}=a_{2}$.
\end{defn}
\begin{example}
$f:\mathbb{Z}\rightarrow\mathbb{Z}$ where $f(x)=x^{3}$ is one-to-one.
\end{example}
\begin{proof}
Suppose that $f(x_{1})=f(x_{2})$ for some $x_{1},x_{2}\in\mathbb{Z}$.
So $(x_{1})^{3}=(x_{2})^{3}\Rightarrow\left((x_{1})^{3}\right)^{\nicefrac{1}{3}}=\left((x_{2})^{3}\right)^{\nicefrac{1}{3}}\Rightarrow x_{1}=x_{2}$
\end{proof}

\begin{xca}
Show that any linear function $f:\mathbb{R}\rightarrow\mathbb{R}$
with a non-zero defined slope is one-to-one. 
\end{xca}
\begin{example}
$f:\mathbb{R}\rightarrow\mathbb{R}$ where $f(x)=|x|$ is not one-to-one.
\end{example}
\begin{itemize}
\item []$|-1|=1=|1|$ but $-1\neq1$ so $f$ is not one-to-one.
\end{itemize}
\begin{prop}
\label{prop:f is 1-1 then f^-1 is fuct}If $f\colon A\rightarrow B$
then $f^{-1}=\{x\colon x\in A\mbox{, }f(x)=b\}$ is a function if
and only if $f$ is one-to-one.
\end{prop}
\begin{xca}
Prove proposition \ref{prop:f is 1-1 then f^-1 is fuct}.
\end{xca}
\begin{defn}
Let $f\colon A\rightarrow B$ be a one-to-one function. We call the
function guarteed to exist by proposition \ref{prop:f is 1-1 then f^-1 is fuct},
$f^{-1}\colon B\rightarrow A$, the \emph{inverse} of $f$.
\end{defn}

\begin{defn}
A function \funct{A}{B} is \emph{onto }if for every $b\in B$ there
is an $a\in A$ such that $f(a)=b$.
\end{defn}
\begin{example}
$f\colon\mathbb{R}\rightarrow\mathbb{R}$ where $f(x)=x^{2}$ is onto.
\end{example}
\begin{proof}
If $b\in\mathbb{R}$ then $\sqrt{b}\in\mathbb{R}$, and $f(\sqrt{b})=\left(\sqrt{b}\right)^{2}=b$.
\end{proof}

\begin{example}
$f\colon\mathbb{Z}\rightarrow\mathbb{R}$ where $f(x)=x^{2}$ is not
onto.
\end{example}
\begin{proof}
$2\in\mathbb{R}$ but $\sqrt{2}\notin\mathbb{Z}$.
\end{proof}
\begin{xca}
Prove that the floor function (see example \ref{ex-floor-funct})
is onto but not one-to-one. 
\end{xca}
\begin{figure}[h]
\hfill%
\begin{minipage}[t]{0.33\textwidth}%
 \begin{tikzpicture}[scale=.6,ele/.style={fill=black,circle,minimum width=.8pt,inner sep=1pt},every fit/.style={ellipse,draw,inner sep=-2pt}]   
	\node[label=center:$\textbf{A}$] (a0) at (0,5) {};
	\node[ele,label=left:$a$] (a1) at (0,4) {};
	\node[ele,label=left:$a$] (a1) at (0,4) {};       
	\node[ele,label=left:$b$] (a2) at (0,3) {};       
	\node[ele,label=left:$c$] (a3) at (0,2) {};   
	\node[ele,label=left:$d$] (a4) at (0,1) {};
	\node[label=center:$\textbf{B}$] (b0) at (4,5) {};	
	\node[ele,,label=right:$1$] (b1) at (4,4) {};   
	\node[ele,,label=right:$2$] (b2) at (4,3) {};   
	\node[ele,,label=right:$3$] (b3) at (4,2) {};   
	\node[ele,,label=right:$4$] (b4) at (4,1) {};
	\node[draw,fit= (a1) (a2) (a3) (a4),minimum width=2cm] {};
	\node[label=center:$f$] 
(f1) at (2,4) {};
	 
	\node[draw,fit= (b1) (b2) (b3) (b4),minimum width=2cm] {};     
	
	\draw[->,thick,shorten <=2pt,shorten >=2pt] (a1) -- (b2);   
	\draw[->,thick,shorten <=2pt,shorten >=2] (a2) -- (b1);   
	\draw[->,thick,shorten <=2pt,shorten >=2] (a3) -- (b3);   
	\draw[->,thick,shorten <=2pt,shorten >=2] (a4) -- (b4);  
\end{tikzpicture}\caption{}{\funct{A}{B} is  one-to-one and onto.}%
\end{minipage}\hfill%
\begin{minipage}[t]{0.33\textwidth}%
\begin{center}
 \begin{tikzpicture}[scale=.6,ele/.style={fill=black,circle,minimum width=.8pt,inner sep=1pt},every fit/.style={ellipse,draw,inner sep=-2pt}]   
	\node[label=center:$\textbf{A}$] (a0) at (0,5) {};
	\node[ele,label=left:$a$] (a1) at (0,4) {};
	\node[ele,label=left:$a$] (a1) at (0,4) {};       
	\node[ele,label=left:$b$] (a2) at (0,3) {};       
	\node[ele,label=left:$c$] (a3) at (0,2) {};   
	\node[ele,label=left:$d$] (a4) at (0,1) {};
	\node[label=center:$\textbf{B}$] (b0) at (4,5) {};	
	\node[ele,,label=right:$1$] (b1) at (4,4) {};   
	\node[ele,,label=right:$2$] (b2) at (4,3) {};   
	\node[ele,,label=right:$3$] (b3) at (4,2) {};   
	\node[ele,,label=right:$4$] (b4) at (4,1) {};
	\node[draw,fit= (a1) (a2) (a3) (a4),minimum width=2cm] {};
	\node[label=center:$f$] 
(f1) at (2,4) {};
	 
	\node[draw,fit= (b1) (b2) (b3) (b4),minimum width=2cm] {};     
	
	\draw[->,thick,shorten <=2pt,shorten >=2pt] (a1) -- (b2);   
	\draw[->,thick,shorten <=2pt,shorten >=2] (a2) -- (b2);   
	\draw[->,thick,shorten <=2pt,shorten >=2] (a3) -- (b4);   
	\draw[->,thick,shorten <=2pt,shorten >=2] (a4) -- (b3);  
\end{tikzpicture}\caption{}{\funct{A}{B} is not one-to-one and is not onto.}
\par\end{center}%
\end{minipage}\hfill%
\begin{minipage}[t]{0.33\textwidth}%
\begin{center}
 \begin{tikzpicture}[scale=.6,ele/.style={fill=black,circle,minimum width=.8pt,inner sep=1pt},every fit/.style={ellipse,draw,inner sep=-2pt}]   
	\node[label=center:$\textbf{A}$] (a0) at (0,5) {};
	\node[ele,label=left:$a$] (a1) at (0,4) {};
	\node[ele,label=left:$a$] (a1) at (0,4) {};       
	\node[ele,label=left:$b$] (a2) at (0,3) {};       
	\node[ele,label=left:$c$] (a3) at (0,2) {};   
	\node[ele,label=left:$d$] (a4) at (0,1) {};
	\node[label=center:$\textbf{B}$] (b0) at (4,5) {};	
	%\node[ele,,label=right:$1$] (b1) at (4,4) {};   
	\node[ele,,label=right:$2$] (b2) at (4,3) {};   
	\node[ele,,label=right:$3$] (b3) at (4,2) {};   
	\node[ele,,label=right:$4$] (b4) at (4,1) {};
	\node[draw,fit= (a1) (a2) (a3) (a4),minimum width=2cm] {};
	\node[label=center:$f$] 
(f1) at (2,4) {};
	 
	\node[draw,fit= (b1) (b2) (b3) (b4),minimum width=2cm] {};     
	
	\draw[->,thick,shorten <=2pt,shorten >=2pt] (a1) -- (b2);   
	\draw[->,thick,shorten <=2pt,shorten >=2] (a2) -- (b2);   
	\draw[->,thick,shorten <=2pt,shorten >=2] (a3) -- (b4);   
	\draw[->,thick,shorten <=2pt,shorten >=2] (a4) -- (b3);  
\end{tikzpicture}\caption{}{\funct{A}{B} is onto but not one-to-one.}
\par\end{center}%
\end{minipage}\hfill
\end{figure}

\begin{defn}
A function that is both one-to-one and onto is called a \emph{bijection.}
\end{defn}
\begin{example}
The function \funct{\mathbb{Z}}{\mathbb{Z}} defined by $f(x)=x+1$
is a bijection.

\begin{proof}
Let $b\in\mathbb{Z}$ then $f(b-1)=b-1+1=b$ so $f$ is onto. Now
suppose that $f(a_{1})=f(a_{2})$. $f(a_{1})=a_{1}+1=a_{2}+1=f(a_{2})$
implies that $a_{1}=a_{2}$. So then $f$ is one-to-one. 
\end{proof}
\end{example}
\begin{thm}
\textbf{\textup{Pigeonhole Principle.}}\textbf{\textup{\emph{ }}}\textup{\emph{If
$n>m$ and $n$ items are placed in $m$ boxes then at least one box
will have two items.}}

Equivalently: If $f\colon A\rightarrow B$ and $|A|>|B|$ then $f$
is not one-to-one.
\end{thm}
\begin{itemize}
\item The pigeonhole principle says that there is at least one element $b\in B$
such that $f(a)=f(a')=b$ for some $a,a'\in A$.

\begin{proof}
Suppose to the contrary that $|A|=n>m=|B|$ and there is a a one-to-one
function $f$ from $A$ to $B$. If $A=\{a_{1},a_{2},\dots,a_{n}\}$
then $f(A)=\{f(a_{1}),f(a_{2}),\dots,f(a_{n})\}$ lists $n$ distinct
elements since by definition $f(a_{j})\neq f(a_{k})$. Then $\{f(a_{1}),f(a_{2}),\dots,f(a_{n})\}=\{b_{1},b_{2},\dots,b_{n}\}$
lists $n$ distinct elements in $B$, but this would mean $n\leq m$
which is a contradiction.
\end{proof}
\end{itemize}
\begin{thm}
\label{thm:order and 1-1 and onto}Let $f$ be a function from $A$
to $B$. The the following hold:

\begin{enumerate}
\item If $f$ is one-to-one then $|B|\geq|A|$.
\item If $f$ is onto then $|A|\geq|B|$.
\item If $f$ is a bijection then $|A|=|B|$.
\end{enumerate}
\end{thm}
Part 1 is true by the pigeonhole principle. Part 2 and 3 are left
as exercises.
\begin{example}
Prove parts 2 and 3 of theorem \ref{thm:order and 1-1 and onto}.
\end{example}
\begin{xca}
Let $f\colon A\rightarrow B$. Prove that if $f$ is onto then $|A|\geq|B|$.
\end{xca}

\begin{xca}
Let $f\colon A\rightarrow B$. Prove that if $f$ is onto then $|A|\geq|B|$.
\end{xca}

\begin{xca}
$\bigstar$ Let the function $f:A\rightarrow B$ be 1-1 and onto,
i.e., $f^{-1}$ exists and $\left|A\right|=\left|B\right|$. Prove
that $\left(f^{-1}\circ f\right):A\rightarrow A$ is the identity
function on $A$ and that $\left(f\circ f^{-1}\right):B\rightarrow B$
is the identity function on $B$.
\end{xca}

\begin{xca}
$\bigstar$ Let $f:A\rightarrow B$ be 1-1 and onto. Prove that $g=f^{-1}$,
if $\left(g\circ f\right):A\rightarrow A$ is the identity function
on $A$ and $\left(f\circ g\right):B\rightarrow B$ is the identity
function on $B$. 
\end{xca}

\begin{xca}
[$\smiley$\textbf{.}]Two sets are said to be \emph{equal }if there
exists a bijective function between them; we write $\left|A\right|=\left|B\right|$.
$B$ is said to have \emph{strictly greater cardinality} than $B$
if there is a one-to-one function from $A$ into $B$, but there does
not exist a bijective function from $A$ into $B$; we write$\left|A\right|<\left|B\right|$.
Prove that $\left|\mathbb{N}\right|<\left|\mathbb{R}\right|$.
\end{xca}

\section{Relations}

In a relation, things are either related or they are not. For example,
the statement ``the woman $x$ is taller than the man $y$'' is
either true or false for any $x$ from the set of woman and $y$ for
the set of men. This statement is either true or false. Relations
are a generalization of the arithmetic relations such as $=$, $>$,
$\leq$, etc.

It is different than the concept of function as it does not determine
particular parings of elements. In the example there may be many woman
taller than any given man. We will investigate binary relations, meaning
we will apply statements to two elements at a time. 
\begin{defn}
Let $A$ and $B$ be sets and $P(a,b)$ be true or false for any $a\in A$
and $b\in B$. Then $\mathcal{R}(A,B,P(a,b))$ is called a \emph{relation.
$P(a,b)$ }is called a \emph{statement}. 
\end{defn}
\begin{itemize}
\item If $P(a,b)$ is true we write $a\mathcal{R}b$ and say that \emph{a
is related to b}.
\item If $P(a,b)$ is false we write $a\cancel{\mathcal{R}}b$ and say that
\emph{a is not related to b.}
\item When $A$, $B$, and $P(a,b)$ are understood, we will write just
$\mathcal{R}$.
\end{itemize}
\begin{example}
Let $A$ be the set of men, $B$ be the set of women, and $P(a,b)$
be the statement ``$x$ is $y$'s husband''

\begin{itemize}
\item []Then Mr.\,Obama$\mathcal{R}$Mrs.\,Obama
\item []Then Mr.\,Obama$\not\mathcal{R}$Mr.\,Rommney
\end{itemize}
\end{example}

\begin{example}
Let $V=\{1,2,3,4\}$ and $U=\{a,b,c\}$ and $\mathcal{R}=\{(1,a),(1,b),(4,c)\}$.

\begin{itemize}
\item []$1\mathcal{R}a$, $1\mathcal{R}b$, and $4\mathcal{R}c$
\item []$1\cancel{\mathcal{R}}c$, $2\cancel{\mathcal{R}}a$, and $3\cancel{\mathcal{R}}c$
(among many)
\end{itemize}
\end{example}
In the above example we did not even need to know the statement; $\mathcal{R}$
told us the relations. We can think of $\mathcal{R}$ as a set of
pairs, but note that it is not a function from $V$ to $U$. Since
then it would map $1\rightarrow a$ and $1\rightarrow b$ -it is not
well defined. 
\begin{xca}
Let $A=\{1,2,3,4\}$, $B=\{5,6\}$, and define $P(a,b)=\mbox{''\emph{a }divides }b\mbox{"}$
for $a\in A$ and $b\in B$. Find $\mathcal{R}(A,B,P)$.
\end{xca}

\begin{example}
Let $A=\{1,2,3\}$ and define $P(a,b)=\mbox{''}a<b\mbox{"}$. Then
$\mathcal{R}=\{(1,2),(1,3),(2,3)\}$ since $1<2$, $1<3$, and $2<3$
are all true.
\end{example}
\begin{itemize}
\item If the two sets in binary relation $\mathcal{R}$ are the same, $\mathcal{R}(A,A,P)$,
then we say that $\mathcal{R}$ \emph{is a relation over $A$.}\footnote{A binary relation over a single set is also called an \emph{endorelation}.}
\end{itemize}
\begin{defn}
Let $\mathcal{R}$ be a relation over a set $A$ then:

\begin{enumerate}
\item If $a\mathcal{R}a$ for all $a\in A$ then $\mathcal{R}$ is \emph{reflexive}.

\begin{enumerate}
\item []equivalently: $(a,a)\in\mathcal{R}$ for all $a\in\mathcal{R}\Rightarrow$$\mathcal{R}$
is \emph{reflexive}.
\end{enumerate}
\item If $a\mathcal{R}b$ implies that $b\mathcal{R}a$ for all $a,b\in A$
then $\mathcal{R}$ is \emph{symmetric}.

\begin{enumerate}
\item []equivalently: $(a,b)\in\mathcal{R}$$\Rightarrow(b,a)\in\mathcal{R}$
for all $a,b\in\mathcal{R}\Rightarrow$$\mathcal{R}$ is \emph{symmetric}.
\end{enumerate}
\item If $a\mathcal{R}c$ whenever $a\mathcal{R}b$ and $b\mathcal{R}c$
for all $a,b,c\in A$ then $\mathcal{R}$ is \emph{transitive}.

\begin{enumerate}
\item []equivalently: $(a,c)\in\mathcal{R}$ whenever $(a,b)\in\mathcal{R}$
and $(b,c)\in\mathcal{R}$ or all $a,b,c\in\mathcal{R}\Rightarrow$$\mathcal{R}$
is \emph{transitive}.
\end{enumerate}
\end{enumerate}
\end{defn}
\begin{example}
Let $\mathcal{R}="|"$ over $\mathbb{Z}$; that is, let $\mathcal{R}$
be defined by the statement ``$x$ divides $y$'' where $x,y\in\mathbb{Z}$.

\begin{enumerate}
\item Is $\mathcal{R}$ reflexive ? Yes. Let $a\in\mathbb{Z}$ if $a|a$
then $a|a$.
\item Is $\mathcal{R}$ symmetric? No. $2|4$ but $4\cancel{|}2$. 
\item Is $\mathcal{R}$ transitive? Yes. Let $a,b,c\in\mathbb{Z}$ such
that if $a|b$ and $b|c$. We can then write $c=bn$ and $b=am$ for
some $m,n\in\mathbb{Z}$. If $b=am$ then $c=(am)n$ which implies
that $a|c$.
\end{enumerate}
\end{example}

\begin{defn}
An \emph{equivalence relation} is a relation which is reflexive, symmetric,
and transitive.
\end{defn}
\begin{example}
Let $\mathcal{R}$ be ``equals'' over $\mathbb{R}$. That is $(a,b)\in\mathcal{R}$
if $a=b$. $\mathcal{R}$ is an equivalence relation. 

\begin{itemize}
\item []Reflexive: $a=a\Rightarrow a=a$.
\item []Symmetric: $a=b\Rightarrow b=a$.
\item []Transitive: suppose that $a=b$ and $b=c$ then $a=b=c\Rightarrow a=c$.
\end{itemize}
\end{example}
\begin{xca}
Each of the following statements, $P(x,y)$, define a relation $R$
in the natural numbers, $\mathbb{N}$, i.e., $\mathcal{R}=(\mathbb{N},\mbox{\ensuremath{\mathbb{N}}},P(x,y))$.
Show whether each of the following relations is reflexive. 

\begin{enumerate}
\item $P(x,y)=$$x$ is less than or equal to $y$
\item $P(x,y)=$$x$ divides $y$
\item $x+y=10$
\item $x$ and $y$ are relatively prime.
\end{enumerate}
\end{xca}

\begin{table}
\noindent Use the following to answer the next two problems. Let $E=\{1,2,3\}$.
Consider the following relations in $E$.

\noindent{}%
\begin{tabular}{ll}
$\mathcal{R}_{1}=\left\{ \left(1,1\right),\left(2,1\right),\left(2,2\right),\left(3,2\right),\left(2,3)\right)\right\} $ & $\mathcal{R}_{4}=\left\{ \left(1,1\right),\left(3,2\right),\left(2,3\right)\right\} $\tabularnewline
$\mathcal{R}_{2}=\left\{ \left(1,1\right)\right\} $ & $\mathcal{R}_{5}=E\times E$\tabularnewline
$\mathcal{R}_{3}=\left\{ \left(1,2\right)\right\} $ & \tabularnewline
\end{tabular}

\caption{Use the above relations to complete}
\end{table}

\begin{xca}
Let $E=\{1,2,3\}$. Consider the following relations over $E$.\\
\begin{tabular}{ll}
$\mathcal{R}_{1}=\left\{ \left(1,1\right),\left(2,1\right),\left(2,2\right),\left(3,2\right),\left(2,3)\right)\right\} $ & $\mathcal{R}_{4}=\left\{ \left(1,1\right),\left(3,2\right),\left(2,3\right)\right\} $\tabularnewline
$\mathcal{R}_{2}=\left\{ \left(1,1\right)\right\} $ & $\mathcal{R}_{5}=E\times E$\tabularnewline
$\mathcal{R}_{3}=\left\{ \left(1,2\right)\right\} $ & \tabularnewline
\end{tabular} \\
Determine whether relation $\mathcal{R}_{1},\mathcal{R}_{2},\mathcal{R}_{3},\mathcal{R}_{4},\mathcal{R}_{5}$
is symmetric.
\end{xca}

\begin{xca}
Let $E=\{1,2,3\}$. Consider the following relations in $E$. Consider
the following relations over $E$.\\
\begin{tabular}{ll}
$\mathcal{R}_{6}=\left\{ \left(1,2\right),\left(2,2\right)\right\} $ & $\mathcal{R}_{9}=\left\{ \left(1,1\right)\right\} $\tabularnewline
$\mathcal{R}_{7}=\left\{ \left(1,2\right),\left(2,3\right),\left(1,3\right),\left(2,1\right),\left(1,1\right)\right\} $ & $\mathcal{R}_{10}=E\times E$\tabularnewline
$\mathcal{R}_{8}=\left\{ \left(1,2\right)\right\} $ & \tabularnewline
\end{tabular}\\
Determine whether each relation $\mathcal{R}_{6},\mathcal{R}_{7},\mathcal{R}_{8},\mathcal{R}_{9},\mathcal{R}_{10}$
is transitive. 
\end{xca}

\begin{xca}
Define the following relation on the real numbers, $\mathcal{R}:x\mathcal{R}y\iff xy\geq0$.
Prove that $\mathcal{R}$ is reflexive and symmetric, but that $\mathcal{R}$
is not an equivalence relation.
\end{xca}

\begin{xca}
Define the following equivalence relation on the set of subsets of
$X$ as follows: $A\mathcal{R}B\iff A\cap B\neq\emptyset$ for $A,B\subseteq X$
Prove or disprove that $\mathcal{R}$ is an equivalence relation. 
\end{xca}

Let $\mathcal{R}$ and $\mathcal{R}'$ be symmetric relations in a
set A. Prove that $\mathcal{R}\cap\mathcal{R}'$ is a symmetric relation
in A. (note that $\mathcal{R}$ and $\mathcal{R}'$ are subsets of
$A\times A$).

\begin{xca}
$\bigstar$ Prove: Let $R$ and $R'$ be symmetric relations in a
set $A$; then $R\cap R'$ is a symmetric relation in $A$.
\end{xca}

\begin{xca}
$\bigstar$ Define $a\mathcal{R}b$ as $a=b\bmod n$ for $a,b\in\mathbb{Z}$
and $n\in\mathbb{Z}^{+}$.\footnote{$a=b\bmod n$ if $n|a-b$ for integers $n,a,b$. For example, $5=1\bmod4$
and $3=3\bmod4$. There is an infinite class of integers which meeting
this criteria; this set is notated $[a]=\{a+kn\colon k\in\mathbb{Z}\}$.
Modular arithemetic is covered in chapter \ref{chap:Congruence}.} $a=b\bmod n$ is an equivalence relation.
\end{xca}

\end{document}


%% LyX 2.4.4 created this file.  For more info, see https://www.lyx.org/.
%% Do not edit unless you really know what you are doing.
\documentclass[oneside,english]{book}
\usepackage[T1]{fontenc}
\usepackage[latin9]{inputenc}
\setcounter{secnumdepth}{3}
\setcounter{tocdepth}{3}
\usepackage{units}
\usepackage{amsmath}
\usepackage{amsthm}
\usepackage{amssymb}
\usepackage{cancel}
\usepackage{geometry}
\geometry{verbose,tmargin=1cm,bmargin=1cm,lmargin=1.5cm,rmargin=1.5cm}
\usepackage{wasysym}

\makeatletter
%%%%%%%%%%%%%%%%%%%%%%%%%%%%%% Textclass specific LaTeX commands.
\theoremstyle{definition}
\newtheorem{defn}{\protect\definitionname}
\theoremstyle{definition}
\newtheorem{example}{\protect\examplename}
\theoremstyle{definition}
\newtheorem{xca}{\protect\exercisename}
\theoremstyle{plain}
\newtheorem{thm}{\protect\theoremname}
\ifx\proof\undefined\
  \newenvironment{proof}[1][\proofname]{\par
    \normalfont\topsep6\p@\@plus6\p@\relax
    \trivlist
    \itemindent\parindent
    \item[\hskip\labelsep
          \scshape
      #1]\ignorespaces
  }{%
    \endtrivlist\@endpefalse
  }
  \providecommand{\proofname}{Proof}
\fi
\theoremstyle{plain}
\newtheorem{lem}{\protect\lemmaname}
\theoremstyle{plain}
\newtheorem{prop}{\protect\propositionname}

%%%%%%%%%%%%%%%%%%%%%%%%%%%%%% User specified LaTeX commands.
\usepackage{hyperref}
\usepackage[usenames,dvipsnames]{xcolor} %adds additional color options
\hypersetup{colorlinks=true, urlcolor=Blue}%Blue

\usepackage{multicol}

\usepackage{tikz}
\usepackage{tikz-3dplot}
\usetikzlibrary{calc,arrows,shapes,backgrounds,matrix,shapes.geometric,fit,trees,positioning}

\usetikzlibrary{patterns}
\usetikzlibrary{plotmarks}
\usepackage{pgfplots}
\pgfplotsset{compat=newest}

\usepackage{calc}
\usepackage[nomessages]{fp}

\usepackage{datetime}

%%%%%commands
\newcommand{\mb}[1]{$\mathbb{#1}$}
\newcommand{\funct}[2]{$f\colon#1\rightarrow#2$}
% Added by lyx2lyx
\setlength{\parskip}{\smallskipamount}
\setlength{\parindent}{0pt}

\makeatother

\usepackage{babel}
\providecommand{\definitionname}{Definition}
\providecommand{\examplename}{Example}
\providecommand{\exercisename}{Exercise}
\providecommand{\lemmaname}{Lemma}
\providecommand{\propositionname}{Proposition}
\providecommand{\theoremname}{Theorem}

\begin{document}

\chapter{Integers and Congruence\label{chap:Congruence}}

\section{Integers}

Much of abstract algebra studies the set of integers, subsets of integers,
or sets which are equivalent to subsets of integers. Thus to move
forward, it will behoove us to step back and investigate the properties
of integers. You have been using many of these properties all of your
life. Taking a rigorous look at these properties will not only enrich
your understanding, but will allow you to apply these fundamental
properties to more general sets. 
\begin{defn}
\textbf{The Well Ordering Principle. }Every nonempty set of positive
integers contains a smallest member. 
\end{defn}
We will treat the well ordering principle as an axiom.\footnote{Meaning we will just assume it to be true for the integers. The well
ordering principle sometimes given as a theorem proved with an induction
(see below). As it turns out, induction and the well ordering principle
are logically equivalent, and we will instead show that induction
follows from the well ordering principle.}
\begin{defn}
We say that a nonzero integer $d$ \emph{divides }an integer $m$
(or say that $d$ is a \emph{divisor of $m$}) if $m$ can be written
as, $m=dn$ for some integer $n$, denoted $d|m$.

\begin{itemize}
\item $d|m\Rightarrow m=dn$ for some integer $n$.
\item $d\cancel{|}m\Rightarrow m\neq dn$ for any integer $n$.
\end{itemize}
\end{defn}
\begin{example}
$3|12$ since $12=3\cdot4$ (which also shows that $4|12$). $3\cancel{|}10$
since there is no integer $n$ such that $3\cdot n=10$ (solving for
$n$, $n=\nicefrac{10}{3}\notin\mathbb{Z}$).
\end{example}
\begin{xca}
Prove that If $a|b$ and $b|c$ then $a|c$.
\end{xca}

\begin{xca}
If $a|b$ and $a|c$ then $a|b\pm c$.
\end{xca}
\begin{defn}
A prime $p$ is an integer greater that has no divisors other than
$1$ and $p$.

\begin{itemize}
\item Equivalently, $p$ is prime if and only if $d|p\Rightarrow d=1\mbox{ or }d=p$.
\end{itemize}
\end{defn}
\begin{example}
$5$ is prime since $1<5$ and $2,3,4\cancel{|}5$.
\end{example}
\begin{xca}
$\bigstar$ Prove that there are infinitely many prime numbers. 
\end{xca}
\begin{thm}
\textbf{\textup{Division Algorithm. }}\textup{\emph{Let $a,b\in\mathbb{Z}$
and $a,b>0$ then there exists a unique $q,r\in\mathbb{Z}$ such that
$a=bq+r$ where $0\leq r<b$.}}\emph{}

\begin{itemize}
\item \emph{$r$ is called the remainder upon dividing $a$ by $b$}. 
\end{itemize}
\end{thm}
\begin{example}
If you were to divide $13$ by $3$, the division algorithm tells
us that $13=3\cdot4+1$; where $a=13$, $b=3$, $q=4$, and $r=1$.
Once we have chosen $a$ and $b$ (or $q$), the theorem tells us
that $a$ can be written uniquely in this form. 
\end{example}
\begin{itemize}
\item Note: we can write $13$ as $13=4\cdot2+5$, $13=4\cdot1+9$, or $13=4\cdot5-7$.
However the remainders are $5,9>4$ in the former cases and, $-7<0$
in the latter case. The assumption that $0\leq r<b$ is what guarantees
uniqueness.
\end{itemize}

\begin{example}
$57=12\cdot4+9$. 

\begin{proof}
There are two parts to prove, existence and uniqueness. First we show
that you can write any integer $a$ as $bq+r$ where $0\leq r<b$.
Let $S=\{a-bq\colon q\in\mathbb{Z}\mbox{ and }a-bq\geq0\}$. First
note that $S$ is nonempty. If $a-b<0$ then let $k=-1$ so that $a-(-1)b>0$.
If $a-b\geq0$ then let $k=1$. 

If $0\in S$ then $a=bk$ which implies that $b|a$ and then we can
let $q=\nicefrac{a}{b}$ (an integer since $b|a$) and $r=0$ so that
$a=b(\nicefrac{a}{b})+0$. 

If $0\notin S$ then we can apply the well ordering principle (since
$S$ is nonempty) to choose the smallest number in $S$, say $r=a-bq$.
Then $a=bq+r$; $r\geq0$ per our definition of $S$ so we need only
show that $r<b$. If $r\geq b$ then $r-b\geq0\Rightarrow a-bq-b=a-b(q+1)$
so that $a-b(a+1)\in S$. But then $a-bq>a-b(q+1)$ and we already
assumed that $a-bq=r$ was the \emph{smallest }number in $S$ --a
contradiction. So $r<b$.

Now we show uniqueness. Let $a=bq+r$ and $a=bq'+r'$ such that $0\leq r<b$
and $0\leq r'<b'$ for integers $q,q',r,r'$. Without loss of generality\footnote{often abbreviated to WLOG, WOLOG or w.l.o.g.; less commonly stated
as without any loss of generality or with no loss of generality. WLOG
means that although a premise of the proof is being narrowed to a
special case, the proof of that special case can be easily applied
to every other case.} we can assume that $r'\geq r$. $bq+r=bq'+r'$ and $bq'+r'-(bq+r)=0\Rightarrow b(q'-q)=r'-r$.
Thus $b|(r'-r)$ and $0\leq r'-r\leq r'<b$. The only way $b$ can
divide a number smaller than itself (and still be an integer) is if
that number is $0$. So then $r'-r=0\Rightarrow r=r'$ which also
shows that $q'=q$.
\end{proof}
\end{example}
\begin{defn}
The \emph{greatest common divisor }(or GCD) of integers $a$ and $b$
is the largest integer that divides both $a$ and $b$; denoted $\gcd(a,b)$.
When $\gcd(a,b)=1$ we say that $a$ and $b$ are \emph{relatively
prime. }
\end{defn}
\begin{example}
$\gcd(10,6)=2$ since $10=2\cdot5$ and $6=2\cdot3$.
\end{example}

\begin{example}
$1100=2^{2}5^{2}11$ and $440=2^{3}5\cdot11$. So $\gcd(1100,440)=2^{2}5\cdot11=220$.
\end{example}
\begin{defn}
The \emph{lowest common multiple }(or LCM) of integers $a$ and $b$
is the smallest integer that is a multiple of both $a$ and $b$.
\end{defn}
\begin{example}
$10=5\cdot2$ and $6=2\cdot3$ so $\mbox{lcm}(10,6)=5\cdot2\cdot3=60$.
Similarly the $\mbox{lcm}(1100,440)=2^{3}5^{2}11=2200$.
\end{example}
\begin{xca}
Determine the following:

\begin{enumerate}
\item $\gcd(2^{4}3^{2}5\cdot7^{2},2\cdot3^{3}7\cdot11)$
\item The lowest common multiple: $\mbox{lcm}(2^{4}3^{2}5\cdot7^{2},2\cdot3^{3}7\cdot11)$.
\end{enumerate}
\end{xca}

\begin{xca}
$\bigstar$ Suppose that $a\neq0$, $b\neq0$, and $\mbox{lcm}(a,b)=m$.
Prove that if $n$ is a common multiple of $a$ and $b$ then $m|n$. 
\end{xca}

\begin{xca}
$\bigstar$ Let $a\neq0$, $b\neq0$, and $d=\gcd\left(a,b\right)$.
Prove that if $d'$ is any common divisor of $a$ and $b$ then $d'|d$.
\end{xca}
The next theorem plays an important role in our study of algebra.
The proof relies on both the well ordering principle and the division
algorithm.
\begin{thm}
\textbf{\emph{\label{thm:GCD-is-a-linear-combo}GCD is a linear combination.
}}For any nonzero integers $a$ and $b$, there exists integers $s$
and $t$ such that $\gcd(a,b)=as+bt$. Furthermore, $\gcd(a,b)$ is
the smallest positive integer that can be written in the form $as+bt$.
\end{thm}
\begin{itemize}
\item In particular, $a$ and $b$ are relatively prime if and only if there
exists integers $s$ and $t$ such that $1=as+bt$.

\begin{proof}
Let $a,b$ be two nonzero integers. Let $S=\{am+bn\colon m,n\in\mathbb{Z}\mbox{ and }am+bn>0\}$.
Note that $S\neq\emptyset$ since if $am+bn\leq0$ then $a(-m)+b(-n)>0$.
By the well ordering principle we can select the smallest number in
$S$, say $d=as+bt$. We want to show that $d=\gcd(a,b)$. 

Using the division algorithm, write $a=dq+r$ where $0\leq r<d$.
If $r>0$ then $r=a-dq=a-(as+bt)q=a-asq-btq=a(1-sq)+b(-tq)\in S$.
But this contradicts the assumption that $d$ is smallest number in
$S$ since $a(1-sq)+b(-tq)<as+bt$. So then $r=0$ implying that $d|a$.
Similarly we can show that $d|b$. This shows that $d$ is a common
divisor of $a$ and $b$. 

Now we show that $d$ is the largest common divisor. Suppose that
$d'$ is another common divisor of $a$ and $b$. So then $a=d'x$
and $b=d'y$ for integers $x$ and $y$. Now $d=as+bt=(d'x)s+('dy)t=d'(xs+yt)$
which implies that $d'|d$. Therefore $d'<d$.
\end{proof}
\end{itemize}
\begin{example}
$\gcd(10,6)=2$ and $2=10(2)+6(-3)$. 
\end{example}

\begin{example}
$\gcd(4,15)=1$. So $4$ and $15$ are relatively prime. $1=4(4)+15(-1)$. 
\end{example}
We will see that theorem \ref{thm:GCD-is-a-linear-combo} is a very
useful theoretical tool. However it does not tell us how to actually
find the linear combinations. Fortunately, we have the Euclidean Algorithm\footnote{From proposition 2 of book seven of the \emph{Elements }written approximately
c.\,300 BC.}.
\begin{example}
\textbf{Euclidean Algorithm. }By repeatedly using the division algorithm
the $\gcd(a,b)$ can always be found by producing a decreasing sequence
of remainders. The algorithm is as follows:
\[
\begin{array}{rclcl}
a & = & bq_{1}+r_{1} & \mbox{ \,\,\,} & 0<r_{1}<b\\
b & = & r_{1}q_{2}+r_{2} &  & 0<r_{2}<r_{1}\\
r_{1} & = & r_{2}q_{3}+r_{3} &  & 0<r_{3}<r_{2}\\
 & \vdots\\
r_{n-3} & = & r_{n-1}q_{n-1}+r_{n-1} &  & 0<r_{n-1}<r_{n-2}\\
r_{n-2} & = & r_{n-1}q_{n}+r_{n} &  & 0<r_{n}<r_{n-1}\\
r_{n-1} & = & r_{n}q_{n+1}+0
\end{array}
\]
By the well ordering principle, this algorithm will eventually result
in a remainder of $0$. Then the last nonzero remainder, $r_{n}$,
is the greatest common divisor.

For example, to find the $\gcd(2520,154)$: 
\[
\begin{array}{rcl}
2520 & = & 154(16)+56\\
154 & = & 56(2)+42\\
56 & = & 42(1)+14\\
42 & = & 14(3)+0
\end{array}
\]
So $\gcd(2520,154)=14$. We can also use this algorithm to find a
linear combination of $2520$ and $154$ that sums to $14$. Using
the above results, we rewrite the remainders as follows: 
\[
\begin{array}{rclcrcl}
56 & = & 42(1)+14 & \mbox{\ensuremath{\Rightarrow}} & 14 & = & 56-42(1)\\
154 & = & 56(2)+42 & \Rightarrow & 14 & = & 56-[154-56(2)]\\
 &  &  &  & 14 & = & (3)56-154\\
2520 & = & 154(16)+56 & \Rightarrow & 14 & = & (3)[2520-154(16)]-154\\
 &  &  &  & 14 & = & 2520(3)+154(-49)
\end{array}
\]
\end{example}
\begin{xca}
Find integers $s$ and $t$ so that $1=7\cdot s+11\cdot t$. Show
that $s$ and $t$ are not unique. 
\end{xca}

\begin{xca}
Let $d=\gcd(a,b)$. If $a=da'$ and $b=db'$, show that $\gcd(a',b')=1$.
\end{xca}

\begin{xca}
$\bigstar$ If $c^{2}=ab$ and $\gcd\left(a,b\right)=1$, prove that
$a$ and $b$ are perfect squares (Hint: By the theorem above, $a$
and $b$ can be written as the product of primes. They are perfect
squares if and only if each prime in this prime factorization has
an even power). 
\end{xca}
\begin{lem}
\textbf{\emph{Euclid's Lemma. }}Let $p$ be a prime and $a$ and $b$
be integers. If $p|ab$ then $p|a$ or $p|b$.\footnote{This lemma first appeared as proposition $30$ in book seven of Euclid's
\emph{Elements}. }

\begin{proof}
Suppose that $p|ab$ but $p\cancel{|}a$. We want to show that $p|b$.
Since $p\cancel{|}a$, $\gcd(p,a)=1$ and by theorem \ref{thm:GCD-is-a-linear-combo}
we can write $1=sp+ta$. Multiplying both sides by $b$ we have $b=bsp+bta$.
Since $p|a$, $a=pn$ for some integer $n$, and 
\begin{eqnarray*}
b & = & bsp+bta\\
\Rightarrow b & = & bsp+btpn\\
\Rightarrow b & = & p(bs+btn)
\end{eqnarray*}
so $p|b$ as needed. 
\end{proof}
\end{lem}
\begin{itemize}
\item Euclid's lemma is not necessarily true when applied to non-primes.
For example, $10|25\cdot2$ but $10\cancel{|}25$ and $10\cancel{|}2$.
\end{itemize}

\section{\textit{\emph{Modular Arithmetic}}}

Modular arithmetic is a system of arithmetic where numbers ``wrap
around'' or ``reset'' after reaching a set value called the \emph{modulus}.
You have probably been using modulus arithmetic all your life without
giving it much thought. One common use is a 12-hour clock. 

If the time is 5:00 then 11 hours later we do not say it is 16:00
(unless you are in the military). Time ``wraps'' around every 12
hours -modulus 12. Using the division algorithm, $16=12(1)+4$, and
we have a remainder of $4$. $16$ ''wraps'' $4$ hours past $12$.
We write this as 
\[
16\equiv4\bmod12
\]
read ``$16$ is congruent/equivalent to $4$ mod $12$''. Adding
$12$ hours to $16$ will wrap around to the same number, $4$, as
will adding $24$, $36$, $48$ etc. 
\begin{eqnarray*}
4 & \equiv & 4\bmod12\\
16 & \equiv & 4\bmod12\\
28 & \equiv & 4\bmod12\\
40 & \equiv & 4\bmod12\\
\vdots\\
4+12(n) & \equiv & 4\bmod12
\end{eqnarray*}
So $4$, $16$, $28$, etc.\,are all members of an infinite set we
call a \emph{class}, denoted $[4]$.
\begin{defn}
Let $a$ and $n$ be integers with $n>0$. The \emph{congruence class
of a modulo n }(denoted $[a]$) is the set of integers that are congruent
to $a$ modulo $n$, i.e., 
\[
[a]=\left\{ b|b\in\mathbb{Z}\mbox{ and }b=a\bmod n\right\} 
\]
\end{defn}
\begin{example}
The congruence class of $3$ modulo $5$ is $[3]=\left\{ \dots,-2,3,8,13,\dots\right\} $.
\end{example}
There are several useful ways to define $a\equiv b\bmod n$. However
when first introduced to the subject it is often easiest to think
of modular arithmetic in terms of the remainder. 
\begin{itemize}
\item If $a=qn+r$ where $q$ is the \emph{quotient }and $r$ is the \emph{remainder
}upon dividing $a$ by the \emph{modulus }$n$, then $a\equiv r\bmod n$. 
\end{itemize}
\begin{example}
If $20$ is divided by $3$ then the remainder is $2$. So $20\equiv2\bmod3$
since $20=(6)3+2$. 
\end{example}

Of course if $20\equiv2\bmod3$ then one would think that $2\equiv20\bmod3$
(as congruence/equivalence should be an equivalence relation), however
dividing $3$ by $20$ is not permissible as any remainder must be
greater than $0$ but less than $20$. For this reason we define $a\equiv b\bmod n$
as follows:
\begin{defn}
$a\equiv b\bmod n$ if and only if $n|(a-b)$.
\end{defn}
\begin{example}
Using the definition above, $20\equiv2\bmod3$ since $3|(20-2)$ and
$3\equiv20\bmod3$ since $3|(2-20)$.
\end{example}

\begin{example}
If If $a=qn+r$ then $r=a-qn\Rightarrow qn=a-r$. So then $n|a-r$
and $a\equiv r\bmod n$.
\end{example}

\begin{example}
Examples of modular arithmetic:
\begin{eqnarray*}
7 & \equiv & 7\bmod13\mbox{\,\,\,\ since }7=(0)13+7\\
20 & \equiv & 7\bmod13\mbox{\,\,\,\ since \ensuremath{20=(1)13+7}}\\
7 & \equiv & 20\bmod13\,\,\,\mbox{since }7=(-1)13+7\\
48 & \equiv & 0\bmod12\,\,\,\mbox{ since }48=(4)12+0\\
75 & \equiv & 25\bmod50\,\,\,\mbox{ since }75=(1)50+25
\end{eqnarray*}
\end{example}
\begin{prop}
If $a\equiv b\bmod n$ and $c\equiv d\bmod n$ for integers $a,b,c,d,n$\footnote{This assumption is only necessary for part 2. While part 1 is true
for all real numbers part may fail for non-integers.} then 

\begin{enumerate}
\item $a+c\equiv(b+d)\bmod n$
\item $ac\equiv bd\bmod n$
\end{enumerate}
\begin{proof}
$a\equiv b\bmod n\Rightarrow n|(a-b)\Rightarrow a-b=k_{1}n$ and $c\equiv d\bmod n\Rightarrow n|(c-d)\Rightarrow k_{2}n$
for some integers $k_{1}$ and $k_{2}$. 

Part 1. $(a+c)-(b+d)=a+c-b-d=(a-b)+(c-d)=k_{1}n+k_{2}n=n(k_{1}+k_{2})$.
Thus $n|[(a+c)-(b+d)]$ which implies that $a+c\equiv(b+d)\bmod n$.

Part 2. $a-b=k_{1}n$ and $c-d=k_{2}n$ implies that $a=k_{1}n+b$
and $c=k_{2}n+d$. 
\begin{eqnarray*}
ac & = & (k_{1}n+b)(k_{2}n+d)\\
ac & = & k_{1}k_{2}n^{2}+k_{1}dn+k_{2}bn+bd\\
ac & = & bd+n(k_{1}k_{2}n+k_{1}d+k_{2}b)\\
ac-bd & = & n(k_{1}k_{2}n+k_{1}d+k_{2}b)
\end{eqnarray*}
So then $n|(ac-bd)$ and $ac\equiv bd\bmod n$.
\end{proof}
\end{prop}
\begin{xca}
Determine $51\mod 13$, $342\mod 85$, and $(82\cdot73)\mod 7$.
\end{xca}

\begin{xca}
Let $n$ be a positive integer greater than $1$ and let $a$ and
$b$ be positive integers. Prove that $a\mod n\equiv b\mod n$ if
and only if $a\equiv b\mod n$. 
\end{xca}

\section{Mathematical Induction}
\begin{example}
Prove that ${\displaystyle \sum_{k=1}^{n}(2k+1)=n^{2}}$. That is,
prove that $1+3+5+\cdots+(2n-1)=n^{2}$.

\begin{proof}
(step 1: base case) Suppose that $n=1$. Then $1=1^{2}=1$. So then
the statement is true for $n=1$.

(step 2: inductive step) Now \emph{assume that the statement is true
for some $n\geq1$}. We want to show that it is then true for $n+1$.
If it is true for some $n\geq1$ then\emph{ ${\displaystyle \sum_{k=1}^{n+1}}(2k+1)={\displaystyle \sum_{k=1}^{n}(2k+1)+2(n+1)-1}$
}(note how we broke it down into a part we knew and the next step
for $n+1$). 

\emph{
\begin{eqnarray*}
{\displaystyle \sum_{k=1}^{n+1}}(2k+1) & = & {\displaystyle \sum_{k=1}^{n}(2k+1)+2(n+1)-1}=\underset{=n^{2}\mbox{ by the assumption}}{\underbrace{\sum_{k=1}^{n}(2k+1)}}+\underset{=2n+1}{\underbrace{2(n+1)-1}}\\
 & = & n^{2}+2n+1\\
 & = & (n+1)^{2}
\end{eqnarray*}
}so the statement is true by induction.
\end{proof}
\end{example}
Once we have shown it is true for $n=1$ in step $1$ then we can
use step $2$ to show that is true for $n=2$. Once we know it is
true for $n=2$ we know it is true for $n=3$ by step $3$ again.
Continuing like this we can show it is true for any $n$. This process
is called mathematical induction or just proof by induction.
\begin{thm}
\textbf{\textup{The Principle of Mathematical Induction. }}\textup{\emph{Let
$S(n)$ be a statement involving a variable $n$, and the following
hold:}}\emph{}

\begin{enumerate}
\item $S(n_{0})$ is true.
\item If $S(n)$ is true for some positive integer $n\geq n_{0}$ then $S(n+1)$
is true.
\end{enumerate}
Then $S(n)$ is true for all positive integers $n$.
\end{thm}
\begin{proof}
Suppose $S(b)$ is true, and $S(n+1)$ is true whenever $S(n)$ is
true, but there is some $n$ for which $S(n)$ is not true. By the
well ordering principle there is a least number , say $k$, such that
$S(k)$ is not true. $k\neq b$ by assumption, but since $k$ is the
smallest number for which $S$ is not true $S(k-1)$ is true. This
is a contradiction since then $S(k)$ is true by assumption. 
\end{proof}
\begin{xca}
Prove that ${\displaystyle \sum_{k=1}^{n}2k=n^{2}+n}$.
\end{xca}

\begin{xca}
Prove that $23^{n}-1$ is divisible by $11$for all positive integers
$n$.
\end{xca}

\begin{xca}
$\bigstar$ The Fibonacci numbers are :$1,1,2,3,5,8,13,21,34,\dots$.
In general, the Fibonacci numbers are defined by $f_{1}=1$, $f_{2}=1$,
and $f_{n+2}=f_{n+1}+f_{n}$ for $n=1,2,3,\dots$. Prove that the
$n^{\mbox{th}}$ Fibonacci number $f_{n}<2^{n}$.
\end{xca}

\section{Additional problems}
\begin{xca}
$\bigstar$ 

Let $a\neq0$, $b\neq0$, and $d=\gcd\left(a,b\right)$. Prove that
if $d'$ is any common divisor of $a$ and $b$ then $d'|d$.
\end{xca}
\begin{thm}
\textbf{\emph{The Fundamental Theorem of Arithmetic. }}Every integer
greater than 1 is either a prime or the unique product of primes. 
\end{thm}

\begin{xca}
$\bigstar$ If $c^{2}=ab$ and $\gcd\left(a,b\right)=1$, prove that
$a$ and $b$ are perfect squares (Hint: By the theorem above, $a$
and $b$ can be written as the product of primes. They are perfect
squares iff. each prime in this prime factorization has an even power). 

\vspace{0cm}

The set of all congruence classes modulo $n$ is denoted $\mathbb{Z}_{n}$
(read ``$Z$ mod $n$''). 

$\left|\mathbb{Z}_{n}\right|=n$. 
\end{xca}

\begin{xca}
$\bigstar$ Suppose that $a\neq0$, $b\neq0$, and $\mbox{lcm}(a,b)=m$.
Prove that if $n$ is a common multiple of $a$ and $b$ then $m|n$. 
\end{xca}

\begin{xca}
[$\smiley$\textbf{.}]Two sets are said to be \emph{equal }if there
exists a bijective function between them; we write $\left|A\right|=\left|B\right|$.
$B$ is said to have \emph{strictly greater cardinality} than $B$
if there is a one-to-one function from $A$ into $B$, but there does
not exist a bijective function from $A$ into $B$; we write$\left|A\right|<\left|B\right|$.
Prove that $\left|\mathbb{N}\right|<\left|\mathbb{R}\right|$.
\end{xca}

\end{document}


%% LyX 2.4.4 created this file.  For more info, see https://www.lyx.org/.
%% Do not edit unless you really know what you are doing.
\documentclass[oneside,english]{book}
\usepackage[T1]{fontenc}
\usepackage[latin9]{inputenc}
\setcounter{secnumdepth}{3}
\setcounter{tocdepth}{3}
\usepackage{units}
\usepackage{amsmath}
\usepackage{amsthm}
\usepackage{amssymb}
\usepackage{cancel}
\usepackage{geometry}
\geometry{verbose,tmargin=1cm,bmargin=1cm,lmargin=1.5cm,rmargin=1.5cm}
\PassOptionsToPackage{normalem}{ulem}
\usepackage{ulem}

\makeatletter

%%%%%%%%%%%%%%%%%%%%%%%%%%%%%% LyX specific LaTeX commands.
%% Because html converters don't know tabularnewline
\providecommand{\tabularnewline}{\\}

%%%%%%%%%%%%%%%%%%%%%%%%%%%%%% Textclass specific LaTeX commands.
\theoremstyle{definition}
\newtheorem{defn}{\protect\definitionname}
\theoremstyle{definition}
\newtheorem{example}{\protect\examplename}
\ifx\proof\undefined\
  \newenvironment{proof}[1][\proofname]{\par
    \normalfont\topsep6\p@\@plus6\p@\relax
    \trivlist
    \itemindent\parindent
    \item[\hskip\labelsep
          \scshape
      #1]\ignorespaces
  }{%
    \endtrivlist\@endpefalse
  }
  \providecommand{\proofname}{Proof}
\fi
\theoremstyle{definition}
\newtheorem{xca}{\protect\exercisename}
\theoremstyle{plain}
\newtheorem{thm}{\protect\theoremname}
\theoremstyle{plain}
\newtheorem{cor}{\protect\corollaryname}
\theoremstyle{plain}
\newtheorem{lem}{\protect\lemmaname}

%%%%%%%%%%%%%%%%%%%%%%%%%%%%%% User specified LaTeX commands.
\usepackage{hyperref}
\usepackage[usenames,dvipsnames]{xcolor} %adds additional color options
\hypersetup{colorlinks=true, urlcolor=Blue}%Blue

\usepackage{multicol}

\usepackage{tikz}
\usepackage{tikz-3dplot}
\usetikzlibrary{calc,arrows,shapes,backgrounds,matrix,shapes.geometric,fit,trees,positioning}

\usetikzlibrary{patterns}
\usetikzlibrary{plotmarks}
\usepackage{pgfplots}
\pgfplotsset{compat=newest}

\usepackage{calc}
\usepackage[nomessages]{fp}

\usepackage{datetime}

%%%%%commands
\newcommand{\mb}[1]{$\mathbb{#1}$}
\newcommand{\funct}[2]{$f\colon#1\rightarrow#2$}
% Added by lyx2lyx
\setlength{\parskip}{\smallskipamount}
\setlength{\parindent}{0pt}

\makeatother

\usepackage{babel}
\providecommand{\corollaryname}{Corollary}
\providecommand{\definitionname}{Definition}
\providecommand{\examplename}{Example}
\providecommand{\exercisename}{Exercise}
\providecommand{\lemmaname}{Lemma}
\providecommand{\theoremname}{Theorem}

\begin{document}

\chapter{Groups}

We naturally think of many sets as having an associated operation.
Consider the familiar set of integers and the operation addition.
Adding integers gives another integer. Furthermore it does not matter
the order that integers are added. There is a special integer, $0$,
that does not affect other numbers when added to them. We can appreciate
that the set of integers under addition has a desirable structure.

We abstract these properties as \emph{group axioms}, and we will see
a highly diverse collection of structures with these same similar
properties. The ubiquity of groups in areas both within and outside
mathematics makes them a central organizing principle of mathematics.
The concept of group developed in the early nineteenth century in
the study of polynomial equations. 

First we formalize the above mentioned ``operation''.

\section{Groups}
\begin{defn}
Let $G$ be a set. A \emph{(binary) operation}\footnote{When parts of statements are given in parenthesis, it usually means
that future statements will leave that part out. So in the from here
on we will simply write ``operation'' as opposed to writing ``binary
operation''.} on the set $G$ is a function $f\colon G\times G\rightarrow G$.
\end{defn}
\begin{itemize}
\item $f$ is just a function sending two elements to a single element.
For example, define $f$ to be addition. $f(3,4)=3+4=7$.
\end{itemize}
\begin{example}
Define $f\colon\mathbb{R}\times\mathbb{R}\rightarrow\mathbb{R}$ as
$x+1$ for $x,y\in\mathbb{R}$.

\begin{itemize}
\item []$f(1,2)=3$ since $1+2=3$ and $f(2,-5)=2+(-5)=-3$ since, $2-5=-3$.
\end{itemize}
\end{example}

\begin{defn}
A \emph{group}, $G$, be a nonempty set with an associated (binary)
operation $\oplus$, that combines two elements to form another element
with the following properties for all $a,b,c\in G$.

\begin{enumerate}
\item Closure: $a\oplus b\in G$.
\item Associativity: $(a\oplus b)\oplus c=a\oplus(b\oplus c)$.
\item Identity: There exist an $0_{G}\in G$ such that $0_{G}\oplus a=a\oplus0_{G}=a$.
\item Inverse: There exists an $a^{-1}\in G$ such that $a\oplus(a^{-1})=0_{G}$.
\end{enumerate}
\end{defn}
\begin{itemize}
\item We use, $\oplus$, because the operation can be many more things than
$+$. Sometimes juxtaposition is used, e.g., $a\oplus b=ab$, in which
case it should not be confused with multiplication. 
\item The a set $G$ with associated operation $\oplus$ is often denoted
$(G,\oplus)$. For example the set of rational numbers under addition
might be written, $(\mathbb{Q},+)$.
\end{itemize}
\begin{example}
The set of rational numbers, $\mathbb{Q}$, under addition is a group.

\begin{proof}
$\mathbb{Q}$ certainly is nonempty; take $\frac{1}{2}\in\mathbb{Q}$
for instance. We can write any two rational numbers with a common
denominator. So WLOG we can write all rational numbers with the same
nonzero integer denominator, say $d$.

Closure: $\frac{x}{d}+\frac{y}{d}=\frac{x+y}{d}$. Since $x+y\in\mathbb{Z}$
it follows that $\frac{x+y}{d}\in\mathbb{Q}$.

Associativity: $\left(\frac{x}{d}+\frac{y}{d}\right)+\frac{z}{d}=\frac{x+y}{d}+\frac{z}{d}=\frac{x+y+z}{d}=\frac{x}{d}+\left(\frac{y}{d}+\frac{z}{d}\right)$.

Identity: $0=\frac{0}{d}$ and $\frac{0}{d}+\frac{x}{d}=\frac{x}{d}=\frac{x}{d}+\frac{0}{d}$

Inverse: $0$ is the denominator so for any $\frac{x}{d}$ we need
a $\frac{y}{d}$ such that $\frac{x}{d}+\frac{y}{d}=0$. Solve this
equation to find $y=-x$ so for any rational number $\frac{x}{d}$
, $\frac{-x}{d}$ is the additive inverse. 

Thus $\mathbb{Q}$ is a group.
\end{proof}
\end{example}
\begin{xca}
\label{ex:Z is a group under +}Show that the set of integers, $\mathbb{Z}$,
is a group under addition. 
\end{xca}

\begin{xca}
Prove that the set of even integers is a group under addition. 
\end{xca}

\begin{xca}
Prove that integers of $n$ multiples for $n\in\mathbb{Z}^{+}$ are
a group under addition.
\end{xca}

\begin{example}
$\left\{ 1,-1,i,-i\}\right\} $ where $i=\sqrt{-1}$ is a group under
multiplication. From the table,

\begin{tabular}{r|rrrr}
$\times$ & $1$ & $-1$ & $i$ & $-i$\tabularnewline
\hline 
$1$ & $1$ & $-1$ & $i$ & $-i$\tabularnewline
$-1$ & $-1$ & $1$ & $-i$ & $i$\tabularnewline
$i$ & $i$ & $-i$ & $-1$ & $1$\tabularnewline
$-i$ & $-i$ & $i$ & $1$ & $-1$\tabularnewline
\end{tabular}we can see that $1$ is the identity, every element has an inverse
($-i\cdot i=1$ for example), and the set is closed under multiplication.
The last property, associativity, follows by associativity of $\mathbb{C}$. 
\end{example}
To show that a set is not a group, it is usually easiest to find an
example where any one of the properties fails.
\begin{example}
$S=\left\{ 1,-1,i\right\} $ is not a group under multiplication. 
\end{example}
\begin{proof}
$-1(i)=-i$ and $-i\notin S$ so it is not closed. Alternatively,
$i\cdot x=1\Rightarrow x=-i$. So the inverse of $i$ is $-i$, $-i\notin S$. 
\end{proof}
\begin{xca}
\label{matrices are a group}Prove that the set of all $2\times2$
matrices with determinant $1$ and entries from $\mathbb{R}$ is a
group under matrix multiplication. 
\end{xca}

\begin{xca}
Show that the following are \emph{not} groups under the given operation.

\begin{enumerate}
\item $S=\left\{ z\colon z\in\mathbb{Z}\mbox{ and }z=\sqrt{5}\right\} $under
addition: $(S,+)$
\item $\mathbb{Q}$, the set of rationals under multiplication: $(\mathbb{Q},\times)$
\item The set of odd integers under addition.
\item $\mathbb{Z}$ under multiplication: $(\mathbb{Z},\times)$
\end{enumerate}
\end{xca}

\begin{xca}
Let $S=\left\{ a+b\sqrt{5}|\,a,b\mbox{ are not both zero and }a,b\in\mathbb{Q}\right\} $.
Prove that $S$ is a group. (Hint: rationalize the denominator to
show that the inverse of $a+b\sqrt{5}$ is in $S$.)
\end{xca}

\begin{xca}
Show that $\left\{ 1,2,3\right\} $ under multiplication modulo $4$
is \emph{not }a group but that $\left\{ 1,2,3,4\right\} $ under multiplication
modulo $5$ is a group. (Example: $2\cdot3=6$ and $6=2\bmod4$ so
$[6]=[2]$ and thus $2\cdot3$ \emph{is} in $\left\{ 1,2,3\right\} $). 
\end{xca}
It is easy to take the next few theorems for granted. Especially the
cancellation law. You probably have been using it for some time without
giving it much thought. However we cannot assume such properties until
it is proven. Imagine trying to solve a simple linear equation without
it.
\begin{thm}
\textbf{\emph{\label{thm:Cancellation-Law.}Cancellation Law. }}Right
and left cancellation hold in a group, i.e., $b\oplus a=c\oplus a\iff b=c$
and $a\oplus b=a\oplus c\iff b=c$ for any elements $a,b,c$ in the
group.

\begin{proof}
Let $G$ be group with elements $a,b,c$. The inverses of the elements,
$a^{-1},b^{-1},c^{-1}$, are also in $G$ as $G$ is a group. First
we show the case for the right. Suppose that $b\oplus a=c\oplus a$
then 
\begin{eqnarray*}
(b\oplus a)\oplus a^{-1} & = & (c\oplus a)\oplus a^{-1}\\
\Rightarrow b\oplus(a\oplus a^{-1}) & = & c\oplus(a\oplus a^{-1})\mbox{ by associativity}\\
\Rightarrow b\oplus0_{G} & = & c\oplus0_{G}\mbox{ since }a\oplus a^{-1}=0_{G}\\
b & = & c
\end{eqnarray*}

Similarly we can show the case for the left. That $a\oplus b=a\oplus c\iff b=c$. 
\end{proof}
\end{thm}

\begin{thm}
\label{thm:identity-is-unique}A group has one and only one identity
element, i.e., a group's identity element is unique.

\begin{proof}
Let $G$ be a group. Since $G$ is a group it has an identity element
$0_{G}$. Now suppose there is some other $e\in G$ such that $e\oplus a=a\oplus e=a$
for any $a\in G$. We need to show that $e=0_{G}$. Since $G$ is
a group $a$ has an inverse $a^{-1}$. 
\begin{eqnarray*}
(0_{G}\oplus a)\oplus a^{-1} & = & (e\oplus a)\oplus a^{-1}\mbox{ since }o_{G}\oplus a=e\oplus a=a\\
0_{G}\oplus(a\oplus a^{-1}) & = & e\oplus(a\oplus a^{-1})\mbox{ by associativity}\\
0_{G}\oplus0_{G} & = & e\oplus0_{G}\mbox{ since }a\oplus a^{-1}=0_{G}\\
0_{G} & = & e
\end{eqnarray*}
\end{proof}
\end{thm}

\begin{thm}
\label{thm:Inverses-are-unique}An element in a group has one and
only one inverse. 

\begin{proof}
Let $0_{G}$ be the unique identity of a group $G$. Suppose there
exists an $a^{-1},x\in G$ such that $a\oplus a^{-1}=a\oplus x=0_{G}$.
Then 
\begin{eqnarray*}
x & = & x\oplus0_{G}\\
 & = & x\oplus(a\oplus a^{-1})\\
 & = & (x\oplus a)\oplus a^{-1}\\
 & = & 0_{G}\oplus a^{-1}=a^{-1}
\end{eqnarray*}
So $x=a^{-1}$ as needed. 
\end{proof}
\end{thm}
\begin{itemize}
\item Note that since the identity of a group is its own inverse, i.e.,
$0_{G}\oplus0_{G}=0_{G}$, theorem \ref{thm:identity-is-unique} can
be proved as a corollary of theorem \ref{thm:Inverses-are-unique}. 
\end{itemize}
\begin{xca}
$\bigstar$ Addition in $\mathbb{Z}_{n}$ is defined as $\left[a\right]+\left[b\right]=\left[a+b\right]$.
Multiplication in $\mathbb{Z}_{n}$ is defined as $\left[a\right]\cdot\left[b\right]=\left[a\cdot b\right]$.
For example, $\mathbb{Z}_{4}=\left\{ \left[0\right],\left[1\right],\left[2\right],\left[3\right]\right\} $
for convenience we write $\mathbb{Z}_{4}=\left\{ 0,1,2,3\right\} $.
$2+3=5$ and $5=1\bmod4$ so $\left[5\right]=\left[1\right]$. Show
that $\mathbb{Z}_{2}\times\mathbb{Z}_{3}$ is a group under modular
addition. 
\end{xca}
\begin{defn}
A group is \emph{abelian}\footnote{Named in honor of the mathematician Niels Henrik Abel\emph{. }Abel
made many outstanding contributions to a variety of fields; though
he died at the age of 26. Abelian groups are also called commutative
groups.}\emph{ }if $a\oplus b=b\oplus a$ for all $a,b$ in the group.
\end{defn}
\begin{example}
$\mathbb{Z}$ is an abelian group under addition.
\end{example}

\begin{xca}
Prove that the set of all $n\times n$ matrices is an abelian group
under \uline{addition}.
\end{xca}

\begin{xca}
Prove that the set of all $2\times2$ matrices with determinant $1$
and entries from $\mathbb{R}$ is a non-abelian group under matrix
multiplication (we already know that it is a group by exercise \ref{matrices are a group}. 
\end{xca}

\begin{xca}
$\bigstar$ Consider a square placed on the Cartesian plane such that
the square's center is at the origin and its corners are on the axes
(like a diamond; see the figure below). Any motion of the square which
keeps the center at the origin and corners on the axes has the same
result as one of eight motions: $r_{0}\equiv\mbox{ rotation of }0^{\circ}$,
$r_{1}\equiv\mbox{ rotation of }90^{\circ}$, $r_{3}\equiv\mbox{ rotation of }180^{\circ}$,
$d_{x}\equiv\mbox{ }x-\mbox{axis reflection}$, $h_{+}\equiv\,y=x\mbox{ reflection}$,
or $h_{-}\equiv y=-x\mbox{ reflection}$. \\
\hspace*{.5cm}\newcommand{\xmax}{1.5}
\newcommand{\xmin}{-1.5}
\newcommand{\xstep}{0} %must be xmin+n
\newcommand{\ymax}{1.5}
\newcommand{\ymin}{-1.5}
\newcommand{\ystep}{0} %must be ymin+n
%%%%%%%
\begin{tikzpicture} 
\begin{axis}[height=3.5cm, 
	width=3.5cm, 
	axis lines=middle,
	x axis line style={-},y axis line style={-}, 
	grid=none,
	%axis line style={-latex}, 
	xmin=\xmin,xmax=\xmax, 
	ymin=\ymin,ymax=\ymax,
	%restrict y to domain=-1:10,
	no markers, 
	xlabel={$$},ylabel={$$},
	ticks=none,
	]
	
	\addplot[very thick,blue] coordinates{(1,0) (0,1)};
	\addplot[very thick,blue] coordinates{(0,1) (-1,0)};
	\addplot[very thick,blue] coordinates{(-1,0) (0,-1)};
	\addplot[very thick,blue] coordinates{(0,-1) (1,0)};

	\node at (axis cs:1.3,.3) {$1$};
	\node at (axis cs:-.3,1.3) {$2$};
	\node at (axis cs:-1.3,-.3) {$3$};
	\node at (axis cs:.3,-1.3) {$4$};
	
\end{axis}
\end{tikzpicture}\\
Show that $D_{4}=\left\{ r_{0},r_{1},r_{3},d_{x},d_{y},h_{+},h_{-}\right\} $
is a non-abelian group under composition of these motions. The set
$D_{4}$ is called the \emph{dihedral group of degree 4}. 
\end{xca}

\begin{thm}
Let $G$ be a group and $a,b\in G$ then $(a\oplus b)^{-1}=b^{-1}\oplus a^{-1}$.

\begin{proof}
The inverse of $(a\oplus b)^{-1}$ is $(a\oplus b)$. Now $(b^{-1}\oplus a^{-1})\oplus(a\oplus b)=b^{-1}\oplus(a^{-1}\oplus a)\oplus b=b^{-1}\oplus0_{G}\oplus b=b^{-1}\oplus b=0_{G}$.
So then $b^{-1}\oplus a^{-1}$ is the inverse of $(a\oplus b)$, and
by theorem \ref{thm:identity-is-unique} 
\end{proof}
\end{thm}
\begin{xca}
Prove that $G$ is an Abelian group if and only is $\left(a\oplus b\right)^{-1}=a^{-1}\oplus b^{-1}$
for all $a,b\in G$. 
\end{xca}

\begin{xca}
Suppose that if $G$ is a group such that if $a,b,c\in G$ and $a\oplus b=c\oplus a$
then $b=c$. Prove that $G$ is an Abelian group, i.e., that $a\oplus b=b\oplus a$
for all $a,b\in G$. Hint: consider the element $a\oplus(b\oplus a)=(a\oplus b)\oplus a$.
\end{xca}

\begin{xca}
$\bigstar$ Prove that if in a group $G$ that $\left(a\oplus b\right)^{2}=a^{2}\oplus b^{2}$
then $G$ is Abelian.
\end{xca}


\section{Subgroups}
\begin{defn}
If a\emph{ }subset $H$ of a group $G$ is a group under the operation
of $G$, then $H$ is a \emph{subgroup }of $G$, denoted $H\leq G$.
\end{defn}
\begin{itemize}
\item Because $H$ is a group, it is implied in this definition that all
subgroups are non-empty.
\end{itemize}
\begin{example}
\label{ex: 2+Z10 is a subgrop of Z10}$H=\left\{ 2,4,6,8,0\right\} $
is a subgroup of the group $\mathbb{Z}_{10}$ under modular addition.

\begin{proof}
From example \ref{ex:Z is a group under +}, we know that $\mathbb{Z}_{10}$
is a group under addition modulo $10$. \textbf{$H\neq\emptyset$}
and $H\subseteq\mathbb{Z}_{10}$. Using the Cayley table,

\begin{center}{\bfseries{}%
\begin{tabular}{c|ccccc}
$\oplus$ & 0 & 2 & 4 & 6 & 8\tabularnewline
\hline 
0 & 0 & 2 & 4 & 6 & 8\tabularnewline
2 & 2 & 4 & 6 & 8 & 0\tabularnewline
4 & 4 & 6 & 8 & 0 & 2\tabularnewline
6 & 6 & 8 & 0 & 2 & 4\tabularnewline
8 & 8 & 0 & 2 & 4 & 6\tabularnewline
\end{tabular}}\end{center}ee can see that it is closed, has identity element $0$
(the same as $\mathbb{Z}_{10}$), and each element has an inverse
($6+4\equiv0\bmod10$ for example). Associativity holds in modular
addition; so $H$ is a group , and hence a subgroup of $\mathbb{Z}_{10}$.
Written $H\leq\mathbb{Z}_{10}$. 
\end{proof}
\end{example}
A subset $H$ of a group $G$ will always have the associativity property
of $G$. So we need only check for closure, existence of an identity
and inverses. The relation between the identity and inverses allows
us check these properties in a single steo with the following theorem.
\begin{thm}
\textbf{\emph{\label{thm:Subgroup-Test.}Subgroup Test. }}Let $G$
be a group and $H$ be a nonempty subset of $G$. Then $H$ is a subgroup
of $G$ if and only if $a\oplus b^{-1}\in H$ for all $a,b\in H$.

\begin{proof}
$(\Rightarrow)$ Suppose that $H\leq G$, and let $a,b\in H$. $H$
is a group every so element has an inverse, i.e., there exists a $b^{-1}\in H$.
Groups are also closed so $a\oplus b^{-1}\in H$.

$(\Leftarrow)$ Suppose that $a\oplus b^{-1}\in H$, and let $a,b\in H$
for some nonempty subset of $G$. 

\begin{enumerate}
\item []Associativity: The operation, $\oplus$, in $H$ is the same as
that of $G$ so $\oplus$ is associative in $H$. 
\item []Identity: By the hypothesis for any $a\in H\Rightarrow a\oplus a^{-1}=0_{G}\in H$.
So $H$ has the same identity element as $G$ (at least one such element
exists since $H\neq\emptyset$).
\item []Inverse: $0_{G}\in H$ (shown above) so for any $a\in H$ we have
that $0\oplus a^{-1}=a^{-1}\in H$.
\item []Closure: We just showed that for and $a,b\in H\Rightarrow a^{-1},b^{-1}\in H$.
$a\oplus b=a\oplus(b^{-1})^{-1}\in H$.
\end{enumerate}
So $H\leq G$.
\end{proof}
\end{thm}
\begin{itemize}
\item Very often, the operation is just addition. Writing $+$ instead of
$\oplus$, the subgroup test is then: 

\begin{itemize}
\item []A nonempty subset $H$ of a group $G$ is a subgroup of $G$ if
and only if $a-b\in H$ for all $a,b\in H$.
\end{itemize}
\end{itemize}
\begin{example}
$H=\left\{ 2,4,6,8,0\right\} $ is a subgroup of the group $\mathbb{Z}_{10}$
under modular addition. 

\begin{proof}
For any $a,b\in H$, $b^{-1}=-b$ since $b-b=0$. Using the subgroup
test, $a-b\in G$ since $\mathbb{Z}_{10}$ is a group. So we need
only show that $2|a-b$. $2|a$ and $2|b\Rightarrow a=2n$ and $b=2m$
for some $m,n\in\mathbb{Z}$, and $2|2(n-m)\Rightarrow2|a-b$. So
then $a-b\in H$. 
\end{proof}
\end{example}

\begin{example}
Show that $\mathbb{Q}$ is a subgroup of the group $\mathbb{R}$ under
addition. 

\begin{proof}
$\mathbb{Q}$ is a nonempty subset of $\mathbb{R}$. Let $\frac{a}{b},\frac{c}{d}\in\mathbb{Q}$.
$\frac{c}{d}^{-1}=-\frac{c}{d}$ since $\frac{c}{d}-\frac{c}{d}=0$.
$\frac{a}{b}-\frac{c}{d}=\frac{ad-bc}{bd}\in\mathbb{Q}$. So then
$\mathbb{Q}\leq\mathbb{R}$ by the subgroup test. 
\end{proof}
\end{example}
\begin{xca}
Let $\left(\mathbb{Q},+\right)$ be the group of rational numbers
under addition. 

\begin{enumerate}
\item Is $\mathbb{Z}$ a subgroup of $\left(\mathbb{Q},+\right)$?
\item Is $\mathbb{Z}^{+}=\left\{ 0,1,2,3,\dots\right\} $ a subgroup of
$\left(\mathbb{Q},+\right)$?
\end{enumerate}
\end{xca}

\begin{xca}
Determine if the following are subgroups.
\end{xca}
\begin{enumerate}
\item Is $\mathbb{Q}$ a subgroup of the group, $\left(\mathbb{C},+\right)$,
the complex numbers under addition?
\item Is $\mathbb{Z}-\left\{ 0\right\} $ a subgroup of the group, $\left(\mathbb{Q}-\left\{ 0\right\} ,\cdot\right)$,
the non-zero rational numbers under multiplication?
\end{enumerate}

\begin{xca}
Prove that the intersection of two subgroups $H$ and $K$ of a group
$G$ is a subgroup.
\end{xca}

\begin{xca}
Let $H_{1},H_{2},\dots$ be subgroups of $G$. Prove that $H_{1}\cap H_{2}\cap\cdots$
is a subgroup of $G$.
\end{xca}

\begin{xca}
Recall that $\mathbb{Z}_{6}=\left\{ 0,1,2,3,4,5\right\} $ is a group
under modular addition. 

\begin{enumerate}
\item Show that $H=\left\{ 0,2,4\right\} $ and $K=\left\{ 0,3\right\} $
are subgroups of $\mathbb{Z}_{6}$.
\item Show that $H\cup K$ is \emph{not }a subgroup of $\mathbb{Z}_{6}$.
\item Is the union of two subgroups of $G$ necessarily a subgroup $G$?
\end{enumerate}
\end{xca}

\begin{xca}
$\bigstar$ Let $\mathbb{Z}_{6}=\left\{ 0,1,2,3,4,5\right\} $ and
$\mathbb{Z}_{4}=\left\{ 0,1,2,3\right\} $. $\mathbb{Z}_{6}\times\mathbb{Z}_{4}$
is a group under modular addition. Show that $H=\left\{ \left(0,0\right),\left(3,0\right),\left(0,2\right),\left(3,2\right)\right\} $
is a subgroup. 
\end{xca}

\begin{xca}
Let $G$ be an abelian group and $H$ be a subgroup of $G$. Prove
that $S(H)=\left\{ x|x\in G\mbox{ and }x\cdot x\in H\right\} $ is
also a subgroup of $G$.
\end{xca}

\begin{xca}
Let $G$ be a group. Prove that $G$ has a subgroup with one element.
This group is called the \emph{trivial group}.
\end{xca}

\section{Finite Groups and Cyclic Groups}

Groups with finitely many elements have interesting arithmetic properties. 
\begin{defn}
The number of elements in a group its \emph{order}, denoted $\left|G\right|$. 
\end{defn}
This is exactly how we defined the order of a set. For example, $\left|\mathbb{Z}_{10}\right|=10$
and from example \ref{ex: 2+Z10 is a subgrop of Z10}, $\left|H\right|=5$. 

\begin{defn}
The \emph{order of an element }$a$ in a group $G$ is the smallest
positive integer $n$ such that $a^{n}=0_{G}$ ($a^{n}=\underset{n\mbox{ times}}{\underbrace{a\oplus a\oplus\dots\oplus a}}).$
\end{defn}
\begin{example}
From example \ref{ex: 2+Z10 is a subgrop of Z10}, $2+2+2+2+2=10\equiv0\bmod10$.
So in $H$, $\left|2\right|=5$. In $\mathbb{Z}_{10}$, $\left|2\right|=5$
and $\left|1\right|=10$. 
\end{example}
\begin{defn}
The set \emph{generated by $a$ }is the set $\left\langle a\right\rangle =\left\{ a^{n}\colon n\in\mathbb{Z}\right\} $.
The element $a$ is called the \emph{generator }of $\left\langle a\right\rangle $.
\end{defn}
\begin{itemize}
\item If $a$ is an element of a group $G$ with $\left|a\right|=n$ then
$a^{k}\in G$ for any integer $k$, and $\left\langle a\right\rangle =\left\{ 0_{G},a,a^{2},\dots,a^{n-1}\right\} $.
\end{itemize}
\begin{example}
Let $2\in\mathbb{Z}_{10}$. $\left\langle 2\right\rangle =\left\{ 2,3,6,8,0\right\} $.
Recall that we have already showed that this is a subgroup of $\mathbb{Z}_{10}$.
Since $\left\langle 2\right\rangle =\left\{ 2,3,6,8,0\right\} $ we
say that $2$ generates the group $\left\{ 2,3,6,8,0\right\} $.
\end{example}

\begin{example}
Find all generators of $\mathbb{Z}_{6}$.

\begin{itemize}
\item []$\left\langle 1\right\rangle =\left\{ 1,2,3,4,5,0\right\} $(note:
$1$ is a generator of $\mathbb{Z}_{n}$ for any $n$) and $\left\langle 5\right\rangle =\left\{ 5,4,3,2,1,0\right\} $
so $\left\langle 1\right\rangle =\left\langle 5\right\rangle =\mathbb{Z}_{6}$.
\end{itemize}
\end{example}
\begin{xca}
Find all generators of $\mathbb{Z}_{8}$ and $\mathbb{Z}_{10}$.
\end{xca}

Some of your favorite sets are generated by a single element. For
example, any integer $n$ can be written as $\pm1\pm1\pm\cdots\pm1=n$.
So $1$ generates $\mathbb{Z}$. Similarly $2$ and $3$ generate
$2\mathbb{Z}$ and $3\mathbb{Z}$ respectively. 
\begin{defn}
Let $G$ be a group. If there is a single element $g\in G$ such that
every element in $G$ is equal to $g^{n}=\underset{n\mbox{ times}}{\underbrace{g\oplus g\oplus\dots\oplus g}}$
for some $n\in\mbox{\ensuremath{\mathbb{Z}}}$ then $G$ is a \emph{cyclic
}group and \emph{$g$ }is called the \emph{generator} of $G$, and
we write $G=\left\langle g\right\rangle $.
\end{defn}
\begin{itemize}
\item $\left\langle g\right\rangle =\left\{ g^{n}\colon n\in\mathbb{Z}\right\} $
\end{itemize}
\begin{xca}
Prove that $\mathbb{Z}$ is a cyclic group. 
\end{xca}

\begin{xca}
Prove that $\mathbb{Z}_{n}$ is a cyclic group under modular addition
for any $n\in\mathbb{Z}$ such that $n>0$. (you can use the fact
that it is a group).
\end{xca}
\begin{thm}
\label{thm:All-cyclic-groups-are-abelian}All cyclic groups are abelian.

\begin{proof}
Let $G=\left\langle a\right\rangle $ be a group. So any $x,y\in G$
can be written as $a^{i}$ and $a^{j}$ respectively. $a^{i}\oplus a^{j}=a^{i+j}=a^{j+i}=a^{i}\oplus a^{j}$.
\end{proof}
\end{thm}

\begin{thm}
\label{thm:<a> is a subgroup of G}If $G$ is a group and $a\in G$
then $\left\langle a\right\rangle \leq G$.

\begin{proof}
$a\in\left\langle a\right\rangle $ so $\left\langle a\right\rangle $
is nonempty, and $G$ is closed so $a^{m}\in G$ for any integer $m$.
So $\left\langle a\right\rangle $ is a nonempty subset of $G$. Now
we use the subgroup test (theorem \ref{thm:Subgroup-Test.}). 

Let $a^{i}$ and $a^{j}$ be two elements in $\left\langle a\right\rangle $.
$a^{i}\oplus a^{-j}=a^{i-j}\in\left\langle a\right\rangle $ since
$i-j\in\mathbb{Z}$. 
\end{proof}
\end{thm}
\begin{itemize}
\item Theorem \ref{thm:All-cyclic-groups-are-abelian} gives us another
way to show that a group is abelian. 
\item Theorem \ref{thm:<a> is a subgroup of G} gives us an easy way to
find abelian subgroups of any group. 
\end{itemize}
\begin{example}
$\left\langle 2\right\rangle ,\left\langle 3\right\rangle ,\dots,\left\langle n\right\rangle $
are all abelian subgroups of $\mathbb{Z}$. 
\end{example}

\begin{thm}
\label{thm:a^i=00003Da^j iff n|i-j}Let $G$ be a group and $a\in G$.
If $\left|a\right|=n$ then $a^{i}=a^{j}$ if and only if $n|i-j$.

\begin{proof}
$(\Leftarrow)$ Let $n|i-j$. Then $i-j=nq$ and 
\[
a^{i-j}=a^{nq}=\left(a^{n}\right)^{q}=(0_{G})^{q}=0_{g}
\]
as needed.

$(\Rightarrow)$ $a^{i}=a^{j}$ implies that $a^{i}\oplus a^{-j}=a^{i-j}=0_{G}$.
By the division algorithm we can write $i-j=qn+r$ for $q,r\in\mathbb{Z}$
and with $0\leq r<n$. 
\[
0_{G}=a^{i-j}=a^{qn+r}=(a^{n})^{q}\oplus a^{r}=(0_{G})^{q}\oplus a^{r}=a^{r}=0_{G}
\]
Thus $r=0$, and $i-j=qn\Rightarrow n|i-j$ as needed.
\end{proof}
\end{thm}
\begin{itemize}
\item An immediate consequence of this theorem is that the order of the
group generated by $a$ is the same as the order of $a$, i.e., $\left|a\right|=\left|\left\langle a\right\rangle \right|$.
\end{itemize}
\begin{example}
Consider $\mathbb{Z}_{30}$ a group under addition modulo $30$. $\mathbb{Z}_{30}$
has identity $0$ and is generated by $1$. $\left|5\right|=6$ since
\[
5+5+5+5+5+5=(6)5\equiv0\bmod30
\]

What is $5^{40}$ $\bmod30$? $5^{4}=4(5)=20$ and $6|40-4$ so then
$5^{40}\equiv20\bmod30$ as well. More generally, $6|x-4\iff x=6q+4$
for some $q\in\mathbb{Z}$. So then $5^{6q+4}\equiv20\bmod30$. 
\end{example}

\begin{cor}
\label{cor: a^k=00003D0 iff n|k}Let $G$ be a group and $a\in G$.
$\left|a\right|=n$ and $a^{k}=0_{G}$ if and only if $n|k$.

\begin{proof}
$a^{k}=a^{k-0}=0_{G}$ we know by the theorem above that $n|k-0$
so $n|k$. The other direction also follows from the theorem.
\end{proof}
\end{cor}
\begin{figure}
\center\begin{tikzpicture}
	\def \n {5} 
	\def \radius {2.5cm} 
	\def \margin {8} % margin in angles, depends on the radius
	\foreach \s in {1,...,\n} {   
		\node[draw,circle] at ({360/\n * (\s - 1)}:
		\radius) {$\s$};   \draw[->, >=latex] ({360/\n * (\s - 1)+\margin}:\radius) arc ({360/\n * (\s - 1)+\margin}:{360/\n * (\s)-\margin}:\radius);
		} 
\end{tikzpicture}

\caption{A graph representing the cyclic group $\mathbb{Z}_{6}$.}
\end{figure}


\section{Permutation Groups}

Permutation groups are functions from a set $A$ back onto $A$.\footnote{Permutations were the first groups studied by mathematicians. In 1832
the French mathematician �variste Galois gave a criterion for the
solvability of some polynomial equations in terms of the symmetry
group of its roots (solutions), and the elements of these groups correspond
to certain permutations of the roots. Arthur Cayley's work on permutations
in 1854 gave the first abstract definition of a finite group. } They can be thought of ways of ``rearranging'' the elements of
a set.
\begin{defn}
A \emph{permutation }of a set $A$ is a function from $A$ to $A$
that is one-to-one and onto (a bijection). A \emph{permutation group
}of a set $A$ is a set of permutations of $A$ that is a group under
function composition.
\end{defn}
\begin{example}
\label{ex: simple permutations }Let $\sigma\colon\left\{ 1,2,3\right\} \rightarrow\left\{ 1,2,3\right\} $
and $\gamma\colon\left\{ 1,2,3\right\} \rightarrow\left\{ 1,2,3\right\} $
be two bijections from $\left\{ 1,2,3\right\} $ to $\left\{ 1,2,3\right\} $
such that 

\begin{tabular}{rlc}
$\sigma(1)=2$ &  & $\gamma(1)=2$\tabularnewline
$\sigma(2)=3$ & and & $\gamma(2)=1$\tabularnewline
$\sigma(3)=1$ &  & $\gamma(3)=3$\tabularnewline
\end{tabular}

Function composition is then well defined, and $\sigma\circ\gamma\colon\{1,2,3\}\rightarrow\{1,2,3\}$
is also a bijective function from $\{1,2,3\}$ to $\{1,2,3\}$. For
example $\gamma(\sigma(1))=\gamma(2)=1$.
\end{example}
There are different ways to represent permutations on a set $A$.
As permutations map every element in $A$ to one and only one element,
the mapping of any element forms a cycle $a\rightarrow b\rightarrow c\rightarrow\dots\rightarrow a$.
This cycle may not include all the elements in $A$. Indeed it may
just be $a\rightarrow a$, but it will always end where it started.
So a map that takes $a\rightarrow b\rightarrow c\rightarrow\dots\rightarrow a$
can be written conveniently as $(a,b,c,\dots,a)$ and $a\rightarrow a$
as simply $a$. Writing permutations in this way is called \emph{cyclic
form}. 
\begin{example}
Consider the permutations (from the example \ref{ex: simple permutations }
above) $\sigma\colon\left\{ 1,2,3\right\} \rightarrow\left\{ 1,2,3\right\} $
and $\gamma\colon\left\{ 1,2,3\right\} \rightarrow\left\{ 1,2,3\right\} $
be two bijections from $\left\{ 1,2,3\right\} $ to $\left\{ 1,2,3\right\} $
such that 

\begin{tabular}{rlc}
$\sigma(1)=2$ &  & $\gamma(1)=2$\tabularnewline
$\sigma(2)=3$ & and & $\gamma(2)=1$\tabularnewline
$\sigma(3)=1$ &  & $\gamma(3)=3$\tabularnewline
\end{tabular}

Graphs representing $\sigma$ and $\gamma$ are shown in figures \ref{fig:A-graph-representing (1,2,3)}
and \ref{fig:Graphs-representing-(1,2)(3)} respectively.
\end{example}
\begin{figure}[h]
\center\begin{tikzpicture}[scale=.75]
	\def \n {3} %number of elements
	\def \x {3} %endnumber
	\def \ftn {$\sigma$} %function
	\def \radius {2cm}
	\def \radiusNode {2.25cm} %places nodes
	\def \label {$(1,2,3)$} %center label
	\def \margin {8} % margin in angles, depends on the radius
	\foreach [count=\i] \s in {1,2,\x} {   
		\node at ({360/\n * (\i - 1)}:
			\radius) {$\s$};
		\node at ({360/\n * (\i - 1) + 360/(2*\n)}:
			\radiusNode) {\ftn};
		\draw[->, >=latex] ({360/\n * (\i - 1)+\margin}:\radius) arc ({360/\n * (\i - 1)+\margin}:{360/\n * (\i)-\margin}:\radius);
		}
	\node {\label};
\end{tikzpicture}

\caption{\label{fig:A-graph-representing (1,2,3)}A graph representing the
permutation of $\sigma=(1,2,3)$ on $\{1,2,3\}$.}
\end{figure}

\begin{figure}[h]
\begin{minipage}[t]{0.5\textwidth}%
\begin{center}
\begin{tikzpicture}[scale=.75]
	\def \n {2} %number of elements
	\def \x {2} %endnumber
	\def \ftn {$\sigma$} %function
	\def \radius {2cm}
	\def \radiusNode {2.25cm} %places nodes
	\def \label {$(1,2)$} %center label
	\def \margin {8} % margin in angles, depends on the radius
	\foreach [count=\i] \s in {1,\x} {   
		\node at ({360/\n * (\i - 1)}:
			\radius) {$\s$};
		\node at ({360/\n * (\i - 1) + 360/(2*\n)}:
			\radiusNode) {\ftn};
		\draw[<-, >=latex] ({360/\n * (\i - 1)+\margin}:\radius) arc ({360/\n * (\i - 1)+\margin}:{360/\n * (\i)-\margin}:\radius);
	}
	\node {\label};
\end{tikzpicture}
\par\end{center}%
\end{minipage}%
\begin{minipage}[t]{0.5\textwidth}%
\begin{center}
\begin{tikzpicture}[scale=.75]
	\def \n {1} %number of elements
	\def \x {3} %endnumber
	\def \ftn {$\sigma$} %function
	\def \radius {2cm}
	\def \radiusNode {2.25cm} %places nodes
	\def \label {$(3)$} %center label
	\def \margin {8} % margin in angles, depends on the radius
	\foreach [count=\i] \s in {\x} {   
		\node at ({360/\n * (\i - 1)}:
			\radius) {$\s$};
		\node at ({360/\n * (\i - 1) + 360/(2*\n)}:
			\radiusNode) {\ftn};
		\draw[<-, >=latex] ({360/\n * (\i - 1)+\margin}:\radius) arc ({360/\n * (\i - 1)+\margin}:{360/\n * (\i)-\margin}:\radius);
		}
	\node {\label};	
\end{tikzpicture}
\par\end{center}%
\end{minipage}

\caption{\label{fig:Graphs-representing-(1,2)(3)}Graphs representing the permutation
of $\gamma=(1,2)(3)$ on $\{1,2,3\}$.}
\end{figure}

\begin{example}
Let $\alpha$ be a permutation on $\{1,2,3,4,5,6\}$ defined as follows: 

\begin{tabular}{rlc}
$\alpha(1)=2$ &  & $\alpha(4)=3$\tabularnewline
$\alpha(2)=1$ & and & $\alpha(5)=5$\tabularnewline
$\alpha(3)=6$ &  & $\alpha(6)=4$\tabularnewline
\end{tabular}

In cycle form $\alpha=(1,2)(3,6,4)(5)$. See the graph in figure 
\begin{figure}[h]
\begin{minipage}[t]{0.33\textwidth}%
\begin{center}
\begin{tikzpicture}[scale=.75]
	\def \n {2} %number of elements
	\def \x {2} %endnumber
	\def \ftn {$\sigma$} %function
	\def \radius {2cm}
	\def \radiusNode {2.25cm} %places nodes
	\def \label {$(1,2)$} %center label
	\def \margin {8} % margin in angles, depends on the radius
	\foreach [count=\i] \s in {1,\x} {   
		\node at ({360/\n * (\i - 1)}:
			\radius) {$\s$};
		\node at ({360/\n * (\i - 1) + 360/(2*\n)}:
			\radiusNode) {\ftn};
		\draw[<-, >=latex] ({360/\n * (\i - 1)+\margin}:\radius) arc ({360/\n * (\i - 1)+\margin}:{360/\n * (\i)-\margin}:\radius);
	}
	\node {\label};
\end{tikzpicture}
\par\end{center}%
\end{minipage}%
\begin{minipage}[t]{0.33\textwidth}%
\begin{center}
\begin{tikzpicture}[scale=.75]
	\def \n {3} %number of elements
	\def \x {4} %endnumber
	\def \ftn {$\sigma$} %function
	\def \radius {2cm}
	\def \radiusNode {2.25cm} %places nodes
	\def \label {$(3,6,4)$} %center label
	\def \margin {8} % margin in angles, depends on the radius
	\foreach [count=\i] \s in {3,6,\x} {   
		\node at ({360/\n * (\i - 1)}:
			\radius) {$\s$};
		\node at ({360/\n * (\i - 1) + 360/(2*\n)}:
			\radiusNode) {\ftn};
		\draw[->, >=latex] ({360/\n * (\i - 1)+\margin}:\radius) arc ({360/\n * (\i - 1)+\margin}:{360/\n * (\i)-\margin}:\radius);
		}
	\node {\label};	
\end{tikzpicture}
\par\end{center}%
\end{minipage}%
\begin{minipage}[t]{0.33\textwidth}%
\begin{center}
\begin{tikzpicture}[scale=.75]
	\def \n {1} %number of elements
	\def \x {5} %endnumber
	\def \ftn {$\sigma$} %function
	\def \radius {2cm}
	\def \radiusNode {2.25cm} %places nodes
	\def \label {$(5)$} %center label
	\def \margin {8} % margin in angles, depends on the radius
	\foreach [count=\i] \s in {\x} {   
		\node at ({360/\n * (\i - 1)}:
			\radius) {$\s$};
		\node at ({360/\n * (\i - 1) + 360/(2*\n)}:
			\radiusNode) {\ftn};
		\draw[<-, >=latex] ({360/\n * (\i - 1)+\margin}:\radius) arc ({360/\n * (\i - 1)+\margin}:{360/\n * (\i)-\margin}:\radius);
		}
	\node {\label};	
\end{tikzpicture}
\par\end{center}%
\end{minipage}

\caption{\label{fig:Graphs-representing-(1,2)(3)-1}Graphs representing the
permutation of $\alpha=(1,2)(3,6,4)(5)$ on $\{1,2,3,4,5,6\}$.}
\end{figure}

Let $\alpha$ be a permutation on $\{1,2,3,4,5,6\}$ defined as follows: 

\begin{tabular}{rlc}
$\alpha(1)=5$ &  & $\alpha(4)=6$\tabularnewline
$\alpha(2)=3$ & and & $\alpha(5)=2$\tabularnewline
$\alpha(3)=1$ &  & $\alpha(6)=4$\tabularnewline
\end{tabular}

Write in $\alpha$ in cycle form.
\end{example}
What we label the elements in $A$ is really arbitrary as a permutation
only maps elements to elements, and any meaning the labels might otherwise
have is irrelevant to this mapping. For example, both the sets $\{bob,sue,bill\}$
or $\{a,b,c\}$ can be represent by the set $\{1,2,3\}$ as they each
have three elements. As such we just as well use the set $\{1,2,\dots,n-1,n\}$.
Each permutations on a set $A=\{1,2,\dots,n-1,n\}$ is an element
in the collection of all permutations on $A$. A collection we will
call the \emph{symmetric group of degree n}, denoted $S_{n}$. 

Of course we first need to show that $S_{n}$ is a group. We have
already seen an example of the group operation we will use, function
composition, in example \ref{ex: simple permutations }. Note that
function operation will always be well defined as the permutations
are all bijections. This operation is can easily be performed using
cycle notation. 
\begin{example}
Let $\sigma=(1,2,3)$ and $\gamma=(1,2)(3)$. $\sigma\circ\gamma$
or just $\sigma\gamma$ as we will write is then 
\[
\begin{array}{rrr}
1\rightarrow2\rightarrow3 & \Rightarrow & (1,3)\\
2\rightarrow1\rightarrow2 & \Rightarrow & (2)\\
3\rightarrow3\rightarrow1 & \Rightarrow & (3,1)
\end{array}
\]
The first and last cycles are equivalent so he have $\sigma\gamma=(1,3)(2)$.

\begin{itemize}
\item Note that by writing a permutation in cycle form it can be inferred
that both these function are bijections of the set $\left\{ 1,2,3\right\} $. 
\end{itemize}
\end{example}

\begin{example}
Let $\alpha=(1,2)(3)(4,5)$ and $\beta=(1,5,3)(2,4)$. 
\begin{eqnarray*}
\alpha\beta & = & (1,4)(2,5,3)
\end{eqnarray*}
\end{example}
\begin{xca}
Let $\alpha=(1,2)(3)(4,5)$ and $\beta=(1,5,3)(2,4)$. Find $\beta\alpha$.
\end{xca}
$S_{n}$ is a nonempty set with a well defined operation. Furthermore
it is easy to see that this operation is closed. 

What would the identity element be? 
\[
e=(1)(2)(3)\dots(n)
\]
$(1)$ means that $1\rightarrow1$. So then $\sigma e=e\sigma$ for
any permutation $\sigma$ in $S_{n}$. 

What about inverses? Permutations are all one-to-one, so each will
have an inverse function. Thus $S_{n}$ is a group.
\begin{example}
Find the inverse of $\beta=(1,5,3)(2,4)$. We need $1\rightarrow5\rightarrow1$.
Also we need $3\rightarrow1\rightarrow3$, and $5\rightarrow3\rightarrow5$
so $(5,1,3)\in\beta^{-1}$. Similarly $2\rightarrow4\rightarrow2$
and $4\rightarrow2\rightarrow4$ implies that $(2,4)\in\beta^{-1}$.
Giving, 
\[
\beta^{-1}=(5,1,3)(2,4)=(1,3,5)(2,4)\mbox{.}
\]
It is easy to check that $\beta\beta^{-1}=(1)(2)(3)(4)(5)=\beta^{-1}\beta$.
\end{example}
\begin{xca}
Find the inverse of $\alpha=(1,2)(3)(4,5)$ and $\alpha\beta=(1,4)(2,5,3)$.
\end{xca}
It is easy to find the number of permutations in $S_{n}$. Let $\sigma\in S_{n}$
and $\sigma(x)=y$ represent the permutation mapping $x\rightarrow y$.
There are $n$ choices for $\sigma(1)$, then $n-1$ choices for $\sigma(2)$,
next there are $n-2$ choices for $\sigma(3)$ etc. So there will
be a total of $n(n-1)(n-2)\cdots1=n!$ permutations in $S_{n}$, i.e.,
$\left|S_{n}\right|=n!$.
\begin{xca}
Find the order of $S_{4}$, $S_{1}$, and $S_{3}$.
\end{xca}

\begin{xca}
Show by example that $S_{3}$ is not abelian. Thus in general symmetric
groups are not abelian. 
\end{xca}

\section{Cosets}

Recall that the elements in $\mathbb{Z}$ each fell into $n$ classes
of $\mathbb{Z}_{n}$ denoted $\left[a\right]$ for an $a\in\mathbb{Z}_{6}$.
For $\mathbb{Z}_{6}$, $2\equiv8\bmod6$ so then $8\in\left[2\right]$.
We can easily find all the elements in $\left[2\right]$ by repeatedly
adding 6 (since we are working with modulo $6$). Similarly we could
find all the elements of each class. 
\begin{eqnarray*}
\left[2\right] & = & \left\{ \dots,-2,0,2,8,14,\dots\right\} \\
\left[3\right] & = & \left\{ \dots,-3,0,3,6,15,\dots\right\} \\
\vdots\\
\left[a\right] & = & \left\{ \dots,a-n,a,a+n,a+2n,\dots\right\} \\
\vdots\\
\left[6\right] & = & \left\{ \dots,-6,0,6,12,\dots\right\} 
\end{eqnarray*}
The elements of each class of $\left[a\right]$ can each be expressed
as $a+kn$ for some $k\in\mathbb{Z}$. It thus would make sense to
write $a+6\mathbb{Z}$ instead of $\left[a\right]$. 

We call sets represented like $3+6\mathbb{Z}=\left\{ \dots,-3,0,3,6,15,\dots\right\} $
\emph{cosets} (note that $6\mathbb{Z}$ is a subgroup of $\mathbb{Z}$).
The idea was first explored by Galois and is key to proving Lagrange's
theorem (the subject of the next section).
\begin{defn}
Let $H$ be a subgroup of $G$ and $a\in G$. 

The set $a\oplus H=\left\{ a\oplus h\colon h\in H\right\} $ is called
the \emph{left coset of H in G containing a}.

The set $H\oplus a=\left\{ h\oplus a\colon h\in H\right\} $ is called
the \emph{right coset of H in G containing a}.
\end{defn}
\begin{example}
\label{ex: Z8  and <2> }Let $G=\mathbb{Z}_{8}$ and $H=\left\langle 2\right\rangle =\left\{ 0,2,4,6\right\} $.
Then 
\begin{eqnarray*}
3\oplus H=3\oplus\left\langle 2\right\rangle  & = & \left\{ 3+0,3+2,3+4,3+6\right\} \\
 & = & \{3,5,7,1\}
\end{eqnarray*}
also since $\mathbb{Z}_{8}$ is abelian, 
\begin{eqnarray*}
H\oplus3=\left\langle 2\right\rangle \oplus3 & = & \left\{ 0+3,2+3,4+3,6+3\right\} \\
 & = & \{3,5,7,1\}
\end{eqnarray*}

The order of $2\in\left\langle 2\right\rangle $, written $\left|2\right|$,
is $4$ since $2+2+2+2=2\cdot4=8\equiv0\bmod8$. $\left|\mathbb{Z}_{8}\right|=8$,
and $2\mid8$, that is $\left|2\right|\mid\left|\mathbb{Z}_{8}\right|$. 
\begin{itemize}
\item We will see that the order of a subgroup (or element) always divides
the order of its group.
\end{itemize}
$4,6,0\in\left\langle 2\right\rangle $ and $4\oplus\left\langle 2\right\rangle =\left\{ 4,6,0,2\right\} =2\oplus\left\langle 2\right\rangle $;
$6\oplus\left\langle 2\right\rangle =\left\{ 4,6,0,2\right\} =2\oplus\left\langle 2\right\rangle $;
and $0\oplus\left\langle 2\right\rangle =\left\langle 2\right\rangle =\left\{ 4,6,0,2\right\} =2\oplus\left\langle 2\right\rangle $\@.
From above $3\notin\{0,2,4,6\}$ and $3\oplus\left\langle 2\right\rangle \neq2\oplus\left\langle 2\right\rangle $. 
\begin{itemize}
\item We will see that cosets are either mutually distinct or equal. As
all elements in $\mathbb{Z}_{8}$ are in $\{0,2,4,6\}$ or $\{1,3,5,7\}$
these are the only \emph{distinct }cosets. Meaning that any other
coset $a\oplus\left\langle 2\right\rangle $ will be congruent to
one of these two cosets. 
\end{itemize}
\end{example}
\begin{xca}
Find all the distinct cosets of $\mathbb{Z}_{9}$.
\end{xca}

\begin{xca}
Consider the group $S_{3}=\{(1),(1,2),(1,3),(2,3)\}$ and its subgroup
$H=\{(1),(1,3)\}$. The distinct left cosets are 

\begin{eqnarray*}
(1)H & = & \left\{ (1)(1),(1)(1,3)\right\} =H\\
(1,2)H & = & \left\{ (1,2)(1),(1,2)(1,3)\right\} =\left\{ (1,2),(1,3,2)\right\} =(1,3,2)H\\
(1,3)H & = & \left\{ (1,3)(1),(1,3)(1,3)\right\} =\left\{ (1,3),(1,3)\right\} =H\\
(2,3)H & = & \left\{ (2,3)(1),(2,3)(1,3)\right\} =\left\{ (2,3),(1,2,3)\right\} =\left(1,2,3\right)H
\end{eqnarray*}
\end{xca}
From these examples/exercises we can see that cosets have some distinct
properties.
\begin{lem}
\label{lem:a is in a+H}Let $H$ be a subgroup of $G$ and $a\in H$.
Then $a\in a\oplus H$ and $a\in H\oplus a$.

\begin{proof}
Let $a\in H$. Since $H$ is a subgroup it has an identity element
$0_{G}$, and $a=a\oplus0_{G}\in a\oplus H$ and $a=0_{G}\oplus a\in a\oplus H$.
\end{proof}
\end{lem}

\begin{lem}
\label{lem:a+H=00003DH iff a in H}Let $H$ be a subgroup of $G$.
If $a\in G$ then $a\oplus H=H$ and $H\oplus a=H$ if and only if
$a\in H$.

\begin{proof}
First we show the result for $a\oplus H=H$ .

$(\Rightarrow)$ Suppose that $a\oplus H=H$. $H$ is a group with
the same identity element as $G$, say $0$. $a\oplus0=a\Rightarrow a\in a\oplus H$
as needed. 

$(\Leftarrow)$ Suppose that $a\in H$ by the lemma above, $a\in a\oplus H$.
As $H$ is close, $a\oplus h\in H$ for any $h\in H$. So then $a\oplus H\subseteq H$. 

Now let $h\in H$. Since $a$ is also in $H$ it follows that $a^{-1}\in H$
and thus $a^{-1}\oplus h\in H$. $h=0\oplus h=(a\oplus a^{-1})\oplus h=a\oplus(a^{-1}\oplus h)\in a\oplus H$.
Thus $H\subseteq a\oplus H$ and it follows that $H=a\oplus H$.

The case for $H\oplus a=H$ follows similarly. 
\end{proof}
\end{lem}
\begin{itemize}
\item The results below are all true for right hand cosets and left hand
cosets. For convenience they are stated for left hand cosets.
\end{itemize}
\begin{thm}
\label{thm:cosets =00003D or disjoint}Let $H$ be a subgroup of $G$.
If $a,b\in G$ then:

\begin{itemize}
\item []$a\oplus H=b\oplus H$ or
\item []$(a\oplus H)\cap(b\oplus H)=\emptyset$. 
\end{itemize}
\begin{proof}
Suppose that the cosets are not disjoint, i.e., $(a\oplus H)\cap(b\oplus H)\neq\emptyset$.
So there is some $x\in(a\oplus H)\cap(b\oplus H)$. Thus there exists
an $h_{1},h_{2}\in H$ such that $x=a\oplus h_{1}$ and $x=b\oplus h_{2}$.
Thus 
\begin{eqnarray*}
a\oplus h_{1} & = & b\oplus h_{2}\\
a & = & (b\oplus h_{2})\oplus h_{1}^{-1}\\
a & = & b\oplus(h_{2}\oplus h_{1}^{-1})
\end{eqnarray*}
Which implies that $a\oplus H=\left(b\oplus(h_{2}\oplus h_{1}^{-1})\right)\oplus H=b\oplus(h_{2}\oplus h_{1}^{-1})\oplus H=b\oplus H$.
As needed. The last equality follows since by closure $h_{2}\oplus h_{1}^{-1}\in H\iff(h_{2}\oplus h_{1}^{-1})\oplus H=H$
by lemma \ref{lem:a+H=00003DH iff a in H}.
\end{proof}
\end{thm}
\begin{cor}
Let $H$ be a subgroup of $G$ and $a,b\in G$. $a\oplus H=b\oplus H$
if and only if $a^{-1}\oplus b\in H$. 

\begin{proof}
$a\oplus H=b\oplus H\iff H=\left(a^{-1}\oplus b\right)\oplus H$ which
implies that $a^{-1}\oplus b\in H$ by the theorem above. 
\end{proof}
\end{cor}
\begin{thm}
\label{thm:|b+H| =00003D |a+H|}Let $H$ be a subgroup of $G$ and
$a,b\in G$. $\left|a\oplus H\right|=\left|b\oplus H\right|$

\begin{proof}
One can show that two sets have the same cardinality by finding a
one-to-one map between them. Consider the map $f\colon a\oplus H\rightarrow b\oplus H$
defined by $f(a\oplus h)=b\oplus h$. By the cancellation property,
theorem \ref{thm:Cancellation-Law.}, $f(a\oplus h_{1})=f(a\oplus h_{2})\Rightarrow b\oplus h_{1}=b\oplus h_{2}\Rightarrow h_{1}=h_{2}$
so $f$ is one-to-one as needed. 
\end{proof}
\end{thm}
\begin{cor}
Let $H$ be a subgroup of $G$ and $a\in G$. $\left|a\oplus H\right|=\left|H\right|$

\begin{proof}
This follows immediately from theorem \ref{thm:|b+H| =00003D |a+H|}
above, since $0_{G}\in G$. So then $\left|a\oplus H\right|=\left|0+H\right|=\left|H\right|$.
\end{proof}
\end{cor}

\begin{thm}
Let $H$ be a subgroup of $G$ and $a\in G$. $a\oplus H$ is a subgroup
of $G$ if and only if $a\in H$.

\begin{proof}
$(\Rightarrow)$ Let $a\oplus H\leq G$. So then $0_{G}\in a\oplus H$.
$0_{G}$ is also in $0_{G}\oplus H=H$ thus $(a\oplus H)\cap H\neq\emptyset$
and it follows that $a\oplus H=H$ by theorem \ref{thm:cosets =00003D or disjoint}. 

$(\Leftarrow)$ Suppose that $a\in H$. This implies that $a\oplus H=H$
by lemma \ref{lem:a+H=00003DH iff a in H}, and $a\oplus H=H\leq G$. 
\end{proof}
\end{thm}

\section{Lagrange's Theorem}

Our study of group theory ends with Lagrange's theorem and its applications.
Lagrange's theorem is the most important result in finite group theory.
It provides a necessary characterization for the order of subgroups
(and consequently elements) of a finite group. Recall example \ref{ex: Z8  and <2> }: 
\begin{example}
$\mathbb{Z}_{8}$ and $\left\langle 2\right\rangle =\left\{ 0,2,4,6\right\} $.

\begin{itemize}
\item $\left\langle 2\right\rangle \leq\mathbb{Z}_{8}$: Since $2\in\mathbb{Z}_{8}$,
$\left\langle 2\right\rangle $ is the cyclic subgroup generated by
$2$. Recall that it is also abelian as all cyclic groups are abelian
(theorem \ref{thm:All-cyclic-groups-are-abelian}).
\item $\left|2\right|=\left|\left\langle 2\right\rangle \right|=4$: Since
$2^{4}=2+2+2+2=8\equiv0\bmod8$. 
\item Note that $4\mid8$: This is no coincidence. We will see that Lagrange's
theorem tell us that this is a necessary condition of all subgroups
and group elements.
\item Recall that are cosets of $\left\langle 2\right\rangle $ are either
distinct or disjoint, e.g., $g_{1}+\left\langle 2\right\rangle =g_{2}+\left\langle 2\right\rangle $
or $\left(g_{1}+\left\langle 2\right\rangle \right)\cap\left(g_{2}+\left\langle 2\right\rangle \right)=\emptyset$.
So the distinct left/right cosets of a subgroup \emph{partition }its
group. This will be important in the proof of Lagrange's theorem.
\end{itemize}
\end{example}
\begin{lem}
\label{lem:G is the union of cosets}If $H$ is a subgroup of a finite
group $G$ then $G$ is the union of the distinct cosets of $H$,
i.e., 
\[
G=a_{1}\oplus H\cup a_{2}\oplus H\cup\cdots\cup a_{k}\oplus H
\]
 for a list of elements $a_{1},a_{2},\dots,a_{k}\in G$ where $a_{i}\oplus H\neq a_{j}\oplus H$
if $i\neq j$ for any $1\leq i,j\leq k\leq$. 

\begin{proof}
Let $a_{1}$ be any element in $G$. If there is no such element $a_{2}\in G$
such that $a_{1}\oplus H\neq a_{2}\oplus H$ then $a_{2}\in a_{1}\oplus H$
for any $a_{2}\in G$. Which implies that $a_{1}\oplus H=G$, and
we are done. 

Otherwise we can either choose another $a_{2}$ such that $a_{1}\oplus H\neq a_{2}\oplus H$.
As above, either we can find another element $a_{3}\in G$ such that
$a_{3}\oplus H$ will be distinct from $a_{1}\oplus H$ and $a_{2}\oplus H$
or $a_{1}\oplus H\cup a_{2}\oplus H=G$. $G$ is finite, so continuing
like this we can compile a list of distinct elements in $G$, $a_{1},a_{2},\dots,a_{k}$
for $k\leq\left|G\right|$, such that $G=a_{1}\oplus H\cup a_{2}\oplus H\cup\cdots\cup a_{k}\oplus H$.
\end{proof}
\end{lem}
While the cosets themselves will be unique the list of elements $a_{1},a_{2},\dots,a_{k}$
is not. However the number of cosets $k$ will always be the same.
This number is important enough to have a name -the \emph{index }of
the subgroup $H$.
\begin{defn}
If $H$ is a subgroup of a finite group $G$ such that $G=a_{1}\oplus H\cup a_{2}\oplus H\cup\cdots\cup a_{k}\oplus H$
then the $k$ is called the \emph{index }of $H$ in $G$; denoted
$\left|G\colon H\right|$. As we will see below, $\left|G\colon H\right|=\nicefrac{\left|G\right|}{\left|H\right|}$
so the index is sometimes notated $\left|G\right|/\left|H\right|$.
\end{defn}
\begin{thm}
\textbf{\emph{Lagrange's Theorem. }}If $H$ is a subgroup of a finite
group $G$ then $\left|H\right|\mid\left|G\right|$. 

\begin{proof}
Let $a_{1}\oplus H,a_{2}\oplus H,\dots,a_{k}\oplus H$ be the distinct
left cosets of $H$ in $G$. By lemma \ref{lem:G is the union of cosets}
above, 
\[
G=a_{1}\oplus H\cup a_{2}\oplus H\cup\cdots\cup a_{k}\oplus H
\]
. Since this union is disjoint we have, 
\[
\left|G\right|=\left|a_{1}\oplus H\right|+\left|a_{2}\oplus H\right|+\cdots+\left|a_{k}\oplus H\right|\mbox{.}
\]
By the corollary to theorem \ref{thm:|b+H| =00003D |a+H|}, each coset
$a_{i}\oplus H$ in this union has exactly $\left|H\right|$ elements,
i.e., $\left|a_{1}\oplus H\right|=\left|H\right|$. So then $\left|G\right|=k\cdot\left|H\right|$.
\end{proof}
\end{thm}
\begin{example}
Consider $\mathbb{Z}_{8}$ and its subgroup $\left\langle 2\right\rangle $.
The order of $\left\langle 2\right\rangle $ is either $1,2,4$, or
$8$ as these are the only divisors of $\left|\mathbb{Z}_{8}\right|=8$
. 
\end{example}

\begin{example}
$H=\left\{ 0,1,2\right\} $ is not a subgroup of $\mathbb{Z}_{8}$
as $\left|H\right|=3$ and $3\nmid8$.
\end{example}

\begin{example}
Let $G$ be a group and $H\leq G$. If $\left|G\right|=35=7\cdot5$.
The order of $H$ is $1,5,7,$ or $35$.
\end{example}

\begin{example}
Let $G$ be a group and $H\leq G$. If $\left|G\right|=40$ and $\left|H\right|=7$.
Then $H\nleq G$.
\end{example}
Note: The converse of Lagrange's theorem is false. Simply because
the order of a subset divides the order of a group \uline{does not
make it a subgroup}.
\begin{example}
(The converse of Lagrange's theorem is false). Let $H=\{1,2,3,4\}\subseteq\mathbb{Z}_{8}$.
$\left|H\right|=4$ and $4\mid8$, but $H\nleq G$ as $0\notin G$.
So it has no identity element.
\end{example}

\begin{cor}
If $G$ is a group and $a\in G$ then $\left|a\right|\mid\left|G\right|$.

\begin{proof}
Recall that $\left|\left\langle a\right\rangle \right|=\left|a\right|$
for any element $a$ in a group $G$, and that $\left\langle a\right\rangle \leq G$.
So as an immediate consequence of Lagrange's theorem $\left|a\right|\mid\left|G\right|$.
\end{proof}
\end{cor}

\begin{cor}
If $G$ is a group with prime order then $G$ is cyclic.

\begin{proof}
As $\left|G\right|>1$, we can choose an element $a\in G$ such that
$a\neq0_{G}$. Since $\left|G\right|=p$ for some prime integer $p$,
$\left|a\right|\mid p$ by Lagrange's theorem. Thus $\left|a\right|=1$
or $p$, but $\left|a\right|\neq1$ since $a\neq0_{G}$. It follows
that $\left|\left\langle a\right\rangle \right|=p$, and since $\left\langle a\right\rangle \subseteq G$
the result follows. 
\end{proof}
\end{cor}

\end{document}


%% LyX 2.4.4 created this file.  For more info, see https://www.lyx.org/.
%% Do not edit unless you really know what you are doing.
\documentclass[oneside,english]{book}
\usepackage[T1]{fontenc}
\usepackage[latin9]{inputenc}
\setcounter{secnumdepth}{3}
\setcounter{tocdepth}{3}
\usepackage{amsmath}
\usepackage{amsthm}
\usepackage{amssymb}
\usepackage{cancel}
\usepackage{geometry}
\geometry{verbose,tmargin=1cm,bmargin=1cm,lmargin=1.5cm,rmargin=1.5cm}

\makeatletter
%%%%%%%%%%%%%%%%%%%%%%%%%%%%%% Textclass specific LaTeX commands.
\theoremstyle{definition}
\newtheorem{defn}{\protect\definitionname}
\theoremstyle{definition}
\newtheorem{example}{\protect\examplename}
\theoremstyle{plain}
\newtheorem{prop}{\protect\propositionname}
\ifx\proof\undefined\
  \newenvironment{proof}[1][\proofname]{\par
    \normalfont\topsep6\p@\@plus6\p@\relax
    \trivlist
    \itemindent\parindent
    \item[\hskip\labelsep
          \scshape
      #1]\ignorespaces
  }{%
    \endtrivlist\@endpefalse
  }
  \providecommand{\proofname}{Proof}
\fi
\theoremstyle{definition}
\newtheorem{xca}{\protect\exercisename}
\theoremstyle{plain}
\newtheorem{thm}{\protect\theoremname}

%%%%%%%%%%%%%%%%%%%%%%%%%%%%%% User specified LaTeX commands.
\usepackage{hyperref}
\usepackage[usenames,dvipsnames]{xcolor} %adds additional color options
\hypersetup{colorlinks=true, urlcolor=Blue}%Blue

\usepackage{multicol}

\usepackage{tikz}
\usepackage{tikz-3dplot}
\usetikzlibrary{calc,arrows,shapes,backgrounds,matrix,shapes.geometric,fit,trees,positioning}

\usetikzlibrary{patterns}
\usetikzlibrary{plotmarks}
\usepackage{pgfplots}
\pgfplotsset{compat=newest}

\usepackage{calc}
\usepackage[nomessages]{fp}

\usepackage{datetime}

%%%%%commands
\newcommand{\mb}[1]{$\mathbb{#1}$}
\newcommand{\funct}[2]{$f\colon#1\rightarrow#2$}
% Added by lyx2lyx
\setlength{\parskip}{\smallskipamount}
\setlength{\parindent}{0pt}

\makeatother

\usepackage{babel}
\providecommand{\definitionname}{Definition}
\providecommand{\examplename}{Example}
\providecommand{\exercisename}{Exercise}
\providecommand{\propositionname}{Proposition}
\providecommand{\theoremname}{Theorem}

\begin{document}

\chapter{\label{chap:Rings}Rings}

\section{Rings and Subrings}

We now continue to generalize the familiar arithmetic operations of
addition and multiplication. Just as with groups, results from the
algebra of arithmetic will be extended to a host of non-numerical
objects, e.g., matrices, functions, polynomials, etc.
\begin{itemize}
\item From here on out we will use the following notations unless stated
otherwise:

\begin{itemize}
\item []$+$ for the ``addition'' operation, e.g., $a+b$ (instead of
$\oplus$).
\item []$\cdot$ or juxtaposition for the ``multiplication'' operation,
e.g., $a\cdot b$ or $ab$ (instead of $\otimes$).
\item []''$-a$'' will denote the additive inverse of an element $a$,
e.g., $-a+a=0$ (instead of ``$a^{-1}$'').
\item []''$a^{-1}$'' will denote the multiplicative inverse of an element
$a$, e.g., $a^{-1}a=1$.
\end{itemize}
\end{itemize}
\begin{defn}
A \emph{ring} is a nonempty set, say $R$, with two binary operations
(usually written as multiplication and addition)\footnote{\emph{But of course the operations are not necessarily addition and
multiplication. $+$ and $\times$; $\oplus$ and $\otimes$; are
among the different notations that may be used in varying texts. }} that satisfies the following:

\begin{enumerate}
\item If $a,b\in R$ then $a+b\in R$. \hfill(closed under addition)
\item $a+(b+c)=(a+b)+c$ \hfill(associative addition)
\item $a+b=b+a$ \hfill(commutative addition)
\item There exists an $0_{R}$ such that $a+0_{R}=0_{R}+a=a$ \hfill(additive
identity)
\item There exists an $-a\in R$ such that $a+(-a)=0_{R}$\\
for every $a\in R$. \hfill(additive inverse)
\item If $a,b\in R$ then $ab\in R$. \hfill(closed under multiplication)
\item $a(bc)=(ab)c$ \hfill(associative multiplication)
\item $a(b+c)=ab+ac$ and $(a+b)c=ac+bc$ \hfill(distribution holds).
\end{enumerate}
\end{defn}
\begin{itemize}
\item Note that parts 1-5 say that a ring is an \emph{abelian group.}\footnote{Commutative groups are also called abelian in honor of Norwegian mathematician
Niels Abel.}\emph{ }
\end{itemize}
\begin{example}
The set of integers, $\mathbb{Z}$, under addition and multiplication
is a ring.
\end{example}

\begin{example}
$\mathbb{Z}_{4}=\{0,1,2,3\}$ ($\mathbb{Z}$ modulo $4$) is also
a ring under addition and multiplication as in $\mathbb{Z}_{n}$ for
any $n\in\mathbb{N}$.
\end{example}

\begin{example}
The set of $2$ by $2$ matrices with real numbers as entries,\\
 $\mathcal{M}_{2}(\mathbb{R})=\left\{ \begin{bmatrix}\begin{array}{rr}
a & b\\
c & d
\end{array}\end{bmatrix}\colon a,b,c,d\in\mathbb{R}\right\} $.
\end{example}
\begin{itemize}
\item []Closure, associativity, and distribution can be shown using the
usual rules of matrix multiplication and addition. $\mathcal{M}_{2}(\mathbb{R})$
is commutative under addition (but not multiplication), the additive
inverse is $\begin{bmatrix}\begin{array}{rr}
-a & -b\\
-c & -d
\end{array}\end{bmatrix}$, the additive identity is $\begin{bmatrix}\begin{array}{rr}
0 & 0\\
0 & 0
\end{array}\end{bmatrix}$, and the multiplicative identity is $\begin{bmatrix}\begin{array}{rr}
1 & 0\\
0 & 1
\end{array}\end{bmatrix}$.
\end{itemize}

\begin{prop}
If $a$ and \textbf{$b$} are two elements in $R$ then $(-a)(-b)=ab$.
\end{prop}
\begin{proof}
By definition $a+(-a)=0$. $(a+(-a))(-b)=0$ implies $a(-b)+(-a)(-b)=0$
by distribution. The additive inverses of $a(-b)$ is unique thus
$a(-b)+(-a)(-b)=0$ implies that $(-a)(-b)=ab$.
\end{proof}
\begin{itemize}
\item In particular for rings like $\mathbb{Z}$ and $\mathbb{R}$, the
product of two negatives is a positive. Note the importance of distribution.
\end{itemize}
\begin{xca}
Show that $\mathbb{Z}_{6}$ is a ring with unity under addition and
multiplication.
\end{xca}

\begin{xca}
Show that the set $\sqrt{2}\mathbb{Z}=\{x\colon x=m+\sqrt{2}n\mbox{ for some \ensuremath{m,n\in\mathbb{Z}}}\}$
is a ring.
\end{xca}

\begin{xca}
Let $R$ and $S$ be rings. Define addition and multiplication on
the Cartesian product $R\times S$ by,
\[
(r,s)+(r',s')=(r+r',s+s')\mbox{ and }(r,s)(r',s')=(rr',ss')
\]
for $r,r'\in R$ and $s,s'\in S$. Prove that $R\times S$ is a ring. 
\end{xca}
\begin{defn}
A \emph{commutative ring }is a ring $R$ such that $ab=ba$ for all
$a,b\in R$.
\end{defn}

\begin{defn}
A ring with \emph{identity}\footnote{The term \emph{unity} is often used in place identity. A ring with
an identity element is said to ``have unity''.}\emph{ }is a ring $R$ with an element $1_{R}$ such that $a1_{R}=1_{R}a$
for all $a\in R$.
\end{defn}
\begin{itemize}
\item It is important to check that a potential $1_{R}$ is both a \emph{right
identity, $a1_{R}=a$}; and a \emph{left identity}, $1_{R}b=b$.
\end{itemize}
\begin{example}
$\mathbb{C}$, $\mathbb{R}$, $\mathbb{Q}$, and $\mathbb{Z}$ are
all commutative rings with identity under the usual addition and multiplication.
\end{example}

\begin{example}
\label{ex-Zn-is-comm-ring}$\mathbb{Z}_{n}=\{0,1,2,\dots,n-1\}$ where
$n\geq2$ is a commutative ring with identity under the usual addition
and multiplication (recall that $n$ is called the \emph{modulus}).
\end{example}
\begin{itemize}
\item To add $x+y\mod n$ find $(x+y)\div n$ and find the remainder, say
$r$, then $(x+y)\mod n=r$. In the same way we can find $(xy)\mod n$.
For example, $(3\cdot7)\mod 8=5$ since $2\cdot8+5=3\cdot7$ (divide
by $8$ and take the remainder). 

\begin{proof}
Since $r<n$, $\mathbb{Z}_{n}$ will be closed under addition and
multiplication. Parts 2, 3, 4, and 7 hold because of the normal rules
of arithmetic. Which also give multiplicative commutativity. $0$
and $1$ are the additive and multiplicative identities respectively.

Which only leaves distribution, part 8, and additive inverse, part
5. Let $a\in\mathbb{Z}_{n}$ then $a<n$ so then $n-a\in\mathbb{Z}_{n}$,
and $a+(n-a)=0$. For distribution, let $a,b,c\in\mathbb{Z}_{n}$.
$a(b+c)\mod n=(ab+ac)\mod n=ab\mod n+ac\mod n$ (this was shown last
semester.
\end{proof}
\end{itemize}

\begin{example}
$E=\{x\colon x=2n\mbox{ for some }n\in\mathbb{Z}\}$, the set of even
integers is a commutative ring. However it does \emph{not }have identity. 
\end{example}
\begin{itemize}
\item Keep in mind that a ring is \emph{not }a group under multiplication.
Elements in a ring do not necessarily have a multiplicative inverse,
and rings do not necessarily have a multiplicative identity. Keep
examples of rings in mind which do not have these properties. 
\end{itemize}
\begin{xca}
Prove that a ring that is cyclic\footnote{Recall that a group $G$ is \emph{cyclic} if for every $x\in G$ there
is some $g\in G$ such that $x=g+g+\cdots+g$. In which case we write
$G=\left\langle g\right\rangle $ and say that $G$ is \emph{generated
}by $g$, e.g., $\mathbb{Z}_{3}=\left\langle 1\right\rangle =\left\langle 2\right\rangle $
since $2+2=1$, $2=2$, and $2+2+2=0$ $\mod 3$.} under addition and commutative (under multiplication).
\end{xca}
\begin{thm}
If a ring has an identity, it is unique. 
\end{thm}
\begin{proof}
Suppose to the contrary that a ring $R$ has two identity elements
$1$ and $1'$. $a1=a$ for all $a\in R$ including $1'\in R$. Thus
$1'\cdot1=1$' and $1\cdot1'=1$ . By the definition of a ring $1'=1\cdot1'=1'\cdot1=1$.
Therefore $1'=1$, and the multiplicative identity is unique.
\end{proof}
\begin{xca}
\label{ex-funct-are-ring}Let $F$ be the set of all functions from
$\mathbb{R}$ to $\mathbb{R}$. Define addition and multiplication
as 
\[
(f+g)(x)=f(x)+g(x)\mbox{ and }(fg)(x)=f(x)g(x).
\]
Show that $F$ is a commutative ring with identity.
\end{xca}

\begin{xca}
If a ring element has a multiplicative inverse, it is unique.
\end{xca}

\section{Subrings}
\begin{defn}
Let $R$ be a ring and $S$ be a subset of $R$.\footnote{Notated: $S\subseteq R$.}
If $S$ is a ring under the same operations as $R$, then $S$ is
a \emph{subring}.
\end{defn}
\begin{example}
$\mathbb{Q}$ is a subring of $\mathbb{R}$, and $\mathbb{Z}$ is
a subring of $\mathbb{Q}$ and $\mathbb{R}$.
\end{example}

\begin{example}
The set of even integers is a subring of $\mathbb{Z}$, but the set
of odd integers is not since it is not closed under addition. 
\end{example}
The following two theorems provide shortcuts when trying to show that
a subset is a subring. Particularity, The second theorem can save
quite a bit of work.
\begin{thm}
Suppose that $R$ is a ring and $S\subseteq R$ such that the following
hold:

\begin{enumerate}
\item If $a,b\in S$ then $a+b\in S$ (closed under addition)
\item If $a,b\in S$ then $ab\in S$ (closed under multiplication)
\item $0_{R}\in S$.
\item If $a\in S$ then there is an $-a\in S$ such that $a+(-a)=0_{R}$.
\end{enumerate}
Then $S$ is a subring of $R$.
\end{thm}
\begin{proof}
Parts 2, 3, 7, and 8 of the Ring definition hold for all elements
in $R$ since it is a ring. So they hold for all elements in $S$.
The remaining requirements 1-4 in the theorem satisfy the remaining
requirements. 
\end{proof}
\begin{itemize}
\item Rings and thus subrings must be nonempty. A requirement which is covered
by 3. Since $0_{R}\in S\Rightarrow$ $S\neq\emptyset$. 
\end{itemize}
\begin{thm}
\label{thm: subring test}A nonempty subset $S$ of a ring $R$ is
a subring if whenever $a,b\in S$:

\begin{enumerate}
\item $ab\in S$ (closed under multiplication)
\item $a-b\in S$ (closed under subtraction)
\end{enumerate}
\end{thm}
\begin{proof}
If $a-b\in S$ for all $a,b\in S$ then $S$ is a subgroup of $R$
(subgroup test), and since $R$ is additively commutative it follows
that $S$ is a abelian subgroup of $R$ under addition (this also
means that $0_{R}$). $S$ is multiplicatively and additively associative;
and $S$ is distributive over addition because these properties hold
in $R$ .
\end{proof}
\begin{xca}
Show that the set $3\mathbb{Z}=\{x\colon x=3n\mbox{ for some \ensuremath{n\in\mathbb{Z}}}\}$
is a subring of $\mathbb{Z}$.
\end{xca}

\begin{xca}
Show that the set $k\mathbb{Z}=\{x\colon x=kn\mbox{ for some \ensuremath{n\in\mathbb{Z}}}\}$
for any integer $k$ is a subring of $\mathbb{Z}$.
\end{xca}

\begin{xca}
Prove that the intersection of any two subrings of a ring $R$ is
also a subring of $R$.
\end{xca}

\end{document}



\end{document}
